\documentclass[10pt,twocolumn,letterpaper]{article}

% ---------- IEEE-conference-style page geometry ----------
\usepackage[top=0.75in,bottom=1in,left=0.625in,right=0.625in,columnsep=0.25in]{geometry}
\usepackage{mathptmx}
\usepackage{courier}
\usepackage[T1]{fontenc}
\usepackage{microtype}
\usepackage{amsmath,amssymb}
\usepackage{graphicx}
\usepackage{booktabs}
\usepackage{array}
\usepackage{multirow}
\usepackage{caption}
\usepackage[nocompress]{cite}
\usepackage{url}
\usepackage[colorlinks=true,linkcolor=blue,citecolor=blue,urlcolor=blue]{hyperref}

\sloppy
\hbadness=10000
\emergencystretch=2em
\setlength{\parindent}{1em}
\setlength{\parskip}{0pt}
\linespread{0.98}

% ---------- IEEE-like section headings ----------
\makeatletter
\renewcommand\thesection{\Roman{section}}
\renewcommand\thesubsection{\Alph{subsection}}
\renewcommand\section{\@startsection{section}{1}{\z@}{1.5ex plus 0.5ex minus 0.2ex}{1ex plus 0.2ex}{\centering\normalfont\footnotesize\scshape}}
\renewcommand\subsection{\@startsection{subsection}{2}{\z@}{1.2ex plus 0.4ex minus 0.2ex}{0.5ex plus 0.1ex}{\normalfont\normalsize\itshape}}
\renewcommand\@seccntformat[1]{\csname the#1\endcsname.\ }
\makeatother

% ---------- captions ----------
\captionsetup[figure]{name={Fig.},labelsep=period,font=footnotesize,justification=justified,singlelinecheck=false,skip=3pt}
\captionsetup[table]{name={TABLE},labelsep=newline,font=footnotesize,labelfont=sc,textfont=sc,justification=centering,skip=3pt}
\renewcommand{\thetable}{\Roman{table}}

% ---------- helpers ----------
\newcommand{\us}{\textunderscore\allowbreak}
\newcommand{\var}[1]{{\small\ttfamily\let\_\us #1}}
\newcommand{\tablesize}{\scriptsize}

\setlength{\textfloatsep}{8pt plus 2pt minus 2pt}
\setlength{\floatsep}{8pt plus 2pt minus 2pt}
\setlength{\intextsep}{8pt plus 2pt minus 2pt}

\hypersetup{pdftitle={Interpretable Risk-Factor Dependency Modeling for Lung Cancer Prediction Using Probabilistic Graphical Models -- Final Report M1-M4},pdfauthor={Priyam Modi, Rushin Modi, Jash Odedara}}

\newcommand{\startms}{\setcounter{section}{0}\setcounter{table}{0}\setcounter{figure}{0}\setcounter{equation}{0}}
\newcommand{\msheader}[3]{\par\vspace{6pt}\begin{center}{\large\bfseries #1}\end{center}\vspace{-2pt}
{\small\bfseries\itshape\noindent\textit{Abstract}---#2\par\vspace{2pt}\noindent\textit{Index Terms}---#3\par}\vspace{2pt}\startms}
\renewcommand{\topfraction}{0.95}\renewcommand{\bottomfraction}{0.9}\renewcommand{\textfraction}{0.05}\renewcommand{\floatpagefraction}{0.8}
\renewcommand{\dbltopfraction}{0.95}\renewcommand{\dblfloatpagefraction}{0.8}
\begin{document}
\twocolumn[{%
\begin{center}
{\LARGE\bfseries Interpretable Risk-Factor Dependency Modeling for\\[2pt] Lung Cancer Prediction Using Probabilistic\\[2pt] Graphical Models}\\[6pt]
{\large Final Report: Milestones 1--4}\\[8pt]
{\large Priyam Modi, Rushin Modi, and Jash Odedara}\\[2pt]
{\itshape Group 1, CSE 516 -- Applied AI}\\
{\itshape School of Engineering and Applied Sciences, Ahmedabad University, Ahmedabad, India}\\[8pt]
\end{center}
\begin{center}\begin{minipage}{0.94\textwidth}\small\bfseries\itshape
\noindent\textit{Abstract}---This report presents the complete project on interpretable risk-factor dependency modeling for lung cancer prediction with probabilistic graphical models, organized as four consecutive milestones: M1 (proposal and scoping), M2 (pipeline, PC structure learning and preliminary results), M3 (results and analysis) and M4 (final evaluation). On the Kaggle Lung Cancer Survey dataset (276 unique records after removing 33 duplicates), a PC-learned Bayesian Network is compared with logistic regression, an RBF-SVM and an MLP. Logistic regression obtains the highest AUROC ($0.938\pm0.009$), the SVM the lowest Brier score ($0.079$), and the BN the lowest ECE ($0.043$) but weak minority-class detection. Exact variable elimination matches brute-force enumeration to $1.11\times10^{-16}$, and preliminary evidence favors exact marginalization over imputation for missing evidence. The PGM's defensible contribution is an inspectable dependency representation and exact posterior inference under partial evidence, not superior classification.
\par\vspace{3pt}\noindent\textit{Index Terms}---Probabilistic graphical models, Bayesian networks, PC algorithm, variable elimination, conditional independence, calibration, lung cancer survey.
\end{minipage}\end{center}\vspace{4pt}}]
\msheader{Milestone 1: Proposal and Scoping Document}%
{Lung cancer risk is shaped by a network of interacting demographic, behavioral, and symptomatic factors rather than by any single variable acting in isolation. Standard classifiers can predict a binary outcome with reasonable accuracy but treat the task as a black box, exposing neither the dependency structure among risk factors nor calibrated, partial-evidence queries. This proposal scopes an Applied AI project that represents the Kaggle Lung Cancer survey dataset as a Probabilistic Graphical Model (PGM)---both a Conditional Independence (CI) graph and a Bayesian Network (BN)---learns its structure and parameters from data, and performs exact inference under partial evidence. We adopt Neural Graphical Models (NGMs) \cite{ngm}, which evaluates on the same dataset, as our State-of-the-Art (SOTA) reference. We define baselines (Logistic Regression, Naive Bayes) and KPIs spanning Representation, Inference, and Learning, and lay out a four-phase experimental plan for the remainder of the semester.}%
{Probabilistic graphical models, Bayesian networks, conditional independence graphs, neural graphical models, lung cancer prediction, interpretability.}
\label{ms1}

\section{Problem Statement}
Lung cancer risk is influenced by an interacting set of demographic, behavioral, and symptomatic factors---smoking history, chronic disease, anxiety, peer pressure, chest pain, coughing, and others---rather than by any one variable in isolation. Standard classifiers (logistic regression, gradient-boosted trees) can predict the binary outcome (lung cancer: yes/no) with reasonable accuracy, but they treat prediction as a black box and do not expose how risk factors and symptoms depend on one another. This project asks: can a probabilistic graphical model that explicitly represents the dependency structure among risk factors and symptoms match the predictive quality of standard classifiers while additionally providing an interpretable map of clinical dependencies and calibrated uncertainty?

\subsection{Motivation and Significance}
\begin{itemize}\itemsep0pt
\item Clinical decision support benefits from models that clinicians can inspect and reason about, not just point predictions.
\item Understanding which symptoms are conditionally dependent given smoking status (or given age) can highlight redundant versus informative screening questions.
\item PGMs naturally support partial-evidence queries (e.g., ``what is $P$(lung cancer) given only 3 of 15 observed symptoms?''), which fixed-input classifiers cannot answer without retraining or imputation.
\end{itemize}

\subsection{Assumptions and Scope}
We commit to one domain (clinical risk prediction) and one principal dataset (Lung Cancer Survey Dataset, Kaggle \cite{kaggle}), as required for the Applied AI track. All 15 symptom/behavioral features are treated as categorical/binary; \var{AGE} is treated as continuous (or binned-ordinal), following the mixed-type handling strategy in \cite{ngm} (Sec.~3.3, Appendix D.1). We assume the dataset's self-reported symptom labels are a reasonable, if noisy, proxy for clinical presentation; label-noise analysis is deferred as a stated limitation. The CDC Birth Cohort Linked Birth--Infant Death dataset \cite{cdc}, also used in \cite{ngm}, is noted as a possible stretch extension but is not the principal dataset here, given its scale ($\sim$4M records, $\sim$500 engineered features) relative to a one-semester scope.

\section{Mapping onto Representation, Inference, and Learning}
\subsection{Representation}
Random variables are the 15 binary/categorical patient attributes (\var{GENDER}, \var{SMOKING}, \var{YELLOW\_FINGERS}, \var{ANXIETY}, \var{PEER\_PRESSURE}, \var{CHRONIC\_DISEASE}, \var{FATIGUE}, \var{ALLERGY}, \var{WHEEZING}, \var{ALCOHOL\_CONSUMING}, \var{COUGHING}, \var{SHORTNESS\_OF\_BREATH}, \var{SWALLOWING\_DIFFICULTY}, \var{CHEST\_PAIN}) plus \var{AGE} and the target \var{LUNG\_CANCER}. Each variable is a node; we recover edges in two complementary ways, mirroring the comparison strategy of the SOTA paper:
\begin{itemize}\itemsep0pt
\item A CI graph via sparse precision-matrix estimation (graphical lasso), analogous to the uGLAD graph used in \cite{ngm}.
\item A Bayesian Network (DAG) via score-based structure learning (hill-climbing with a BIC score), analogous to the bnlearn Tabu-search graph in \cite{ngm}.
\end{itemize}
Undirected CI graphs expose symmetric dependence (which symptoms co-occur regardless of ordering), while the DAG exposes a candidate generative ordering (e.g., \var{SMOKING}$\to$\var{CHRONIC\_DISEASE}$\to$\var{LUNG\_CANCER}). Comparing both, as in \cite{ngm}, tests how sensitive downstream inference is to graph choice---itself one of our KPIs (Section~\ref{sec:m1kpi}).

\subsection{Inference}
Target queries are (i) $P(\text{\var{LUNG\_CANCER}}=\text{Yes}\mid\text{observed symptom subset})$, the primary clinical query under partial evidence, and (ii) $P(\text{\var{SMOKING}}\mid\text{remaining features})$, a secondary diagnostic query used in \cite{ngm} (Table~4) to test whether the model robustly reconstructs any variable, not only the labeled outcome. We use exact inference (variable elimination / junction tree) on the learned Bayesian Network for point and marginal queries. As a stretch goal, we plan to implement the gradient-based MAP inference procedure of \cite{ngm} (Algorithm~2) on a small NGM-style network sharing the same graph, to compare inference behavior between a classical BN and a neural graphical model.

\subsection{Learning}
The model is obtained through both structure and parameter learning: hill-climbing with BIC score (Bayesian Network) and graphical-lasso sparse precision estimation (CI graph) for structure; Maximum Likelihood / Bayesian (Dirichlet-prior) estimation of conditional probability tables given the learned DAG for parameters. As an extension aligned with the SOTA paper, we plan to fit a small MLP-based Neural Graphical Model (\cite{ngm}, Eq.~1--2) on the recovered graph, jointly optimizing regression fit and a soft structure-matching penalty, to test whether a richer functional form improves calibration without sacrificing the interpretability of a fixed graph.

\section{State-of-the-Art Position}
Table~\ref{tab:m1sota} summarizes our SOTA reference, Neural Graphical Models \cite{ngm} (Microsoft Research, arXiv preprint 2210.00453, v4).

\begin{table}[!htbp]
\centering
\caption{SOTA position: Neural Graphical Models \cite{ngm}.}
\label{tab:m1sota}
\tablesize
\renewcommand{\arraystretch}{1.15}
\begin{tabular}{@{}>{\bfseries}p{0.2\columnwidth}p{0.74\columnwidth}@{}}
\toprule
Field & Details \\
\midrule
Problem & Represent arbitrary (directed / undirected / mixed-edge, mixed-datatype) feature-dependency graphs with one neural-network-based PGM, supporting efficient learning, inference and sampling. \\
Representation & Each feature is a function of its graph neighbors (or Markov blanket for DAGs), re-parameterized as an MLP whose layer-weight-magnitude product must respect the zero pattern of $G$. \\
Learning & Joint multi-task regression minimizing reconstruction error of $X$ subject to a soft Hadamard-product penalty enforcing near-zero weight paths where $G$ has none; proximal initialization warm-starts the weights. \\
Inference & Message passing, and a gradient-based MAP update that freezes trained weights and updates only unobserved input values until reconstruction matches observed inputs. \\
Evaluation & Synthetic GGMs (graph-recovery AUROC/AUPR); a $\sim$300-patient lung-cancer dataset (5-fold CV vs.\ logistic regression); a 4M-record CDC infant-mortality dataset (vs.\ logistic regression, Bayesian Networks, EBMs). \\
Findings & NGMs match or exceed Bayesian Networks and are competitive with EBMs using a single trained model for all variables, and are notably strong on rare-class prediction. \\
Limitations & Optimization is sensitive to $\lambda$ initialization; scaling to very high-dimensional inputs is an open challenge; text/image inputs are proposed but not empirically evaluated. \\
Relation to our project & \cite{ngm} evaluates on the same Lung Cancer dataset we use, reporting both a Bayesian Network baseline and an NGM model (Table~4). This gives a direct, paper-verified reference point: we reproduce the classical-PGM side ourselves and use their reported NGM methodology/numbers as the SOTA target, with an optional NGM-style extension later. \\
\bottomrule
\end{tabular}
\end{table}

\section{Baseline and Initial KPIs}
\label{sec:m1kpi}
Success is defined jointly across the three PGM verticals, not by predictive accuracy alone (Table~\ref{tab:m1kpi}). Our baseline is Logistic Regression, matching the exact protocol used in \cite{ngm} (Table~4) for this dataset, plus Gaussian Naive Bayes as a second, PGM-flavored baseline assuming full feature independence given the class.

\begin{table}[!htbp]
\centering
\caption{Baselines and KPIs by PGM vertical.}
\label{tab:m1kpi}
\tablesize
\renewcommand{\arraystretch}{1.15}
\begin{tabular}{@{}p{0.2\columnwidth}p{0.45\columnwidth}p{0.27\columnwidth}@{}}
\toprule
Vertical & KPI & Measured against \\
\midrule
Representation & Graph density recovered by CI-graph vs.\ BN; agreement with clinically expected edges (e.g., \var{SMOKING}--\var{COUGHING}) & Cross-method agreement; domain sanity check \\
Representation & Edge stability under bootstrap resampling & Variance across 20 resampled fits \\
Inference & Log-likelihood of held-out patients & Naive Bayes \\
Inference & Calibration: Brier score, ECE & Logistic Regression \\
Inference & Query latency under partial evidence & Baseline classifier inference time \\
Learning & 5-fold CV AUROC / AUPR / accuracy & Logistic Regression \cite{ngm}, Naive Bayes \\
Learning & Sample efficiency (20\%--100\% of data) & Logistic Regression learning curve \\
\bottomrule
\end{tabular}
\end{table}

\section{Experimental Plan, Responsibilities, and Timeline}
\subsection{Experimental Plan}
\begin{itemize}\itemsep0pt
\item \textbf{Phase 1 (M1):} data acquisition and cleaning, exploratory analysis, CI-graph and BN structure recovery, baseline implementation.
\item \textbf{Phase 2 (M2):} parameter learning on recovered graphs, exact inference implementation, KPI evaluation (Section~\ref{sec:m1kpi}), initial comparison against \cite{ngm}.
\item \textbf{Phase 3 (M3):} sensitivity analysis (CI graph vs.\ BN graph, as in \cite{ngm} Fig.~9), calibration analysis, and, time permitting, a small NGM-style neural-view extension trained on the same graph.
\item \textbf{Phase 4 (Final):} consolidation, ablations, writeup, and reflection against the SOTA position.
\end{itemize}

\subsection{Team Responsibilities}
Priyam Modi owns the core PGM pipeline end-to-end---representation (both graph-recovery methods), structure and parameter learning, inference, SOTA analysis, and integration of the report and video---while Rushin Modi and Jash Odedara own data engineering and the baseline/KPI evaluation that the PGM pipeline is compared against (Table~\ref{tab:m1team}).

\begin{table}[!htbp]
\centering
\caption{Team responsibilities (Milestone 1).}
\label{tab:m1team}
\tablesize
\renewcommand{\arraystretch}{1.15}
\begin{tabular}{@{}>{\bfseries}p{0.2\columnwidth}p{0.72\columnwidth}@{}}
\toprule
Member & Primary responsibility (M1) \\
\midrule
Priyam Modi & PGM mapping (Sec.~II): Representation, Inference, and Learning design. SOTA analysis of \cite{ngm} (Sec.~III); KPI design (Sec.~IV). \\
Rushin Modi & Problem framing (Sec.~I); dataset and preprocessing; baseline and experimental plan (Sec.~IV, Sec.~V). \\
Jash Odedara & KPI definitions (Sec.~IV); Naive Bayes baseline; repository structure (Code/, Results/). \\
\bottomrule
\end{tabular}
\end{table}

\section{Conclusion}
This proposal commits our Applied AI project to modeling lung-cancer risk-factor dependencies with a Probabilistic Graphical Model, combining a learned Conditional Independence graph and a Bayesian Network for representation, exact inference under partial evidence, and structure-plus-parameter learning. Using Neural Graphical Models \cite{ngm} as our SOTA reference gives us a directly comparable evaluation on the same dataset, and our KPIs (Section~\ref{sec:m1kpi}) ensure success is judged across representation, inference, and learning rather than accuracy alone.

\msheader{Milestone 2}%
{This report describes an experimental pipeline for learning conditional-dependency structure in a small lung-cancer survey dataset and evaluating probabilistic inference alongside conventional classifiers. The structural-learning method is the PC algorithm: it begins with a complete undirected graph, removes edges using conditional-independence tests while storing separating sets, and then orients only those edges justified by collider and propagation rules. The resulting partially directed acyclic graph (CPDAG) is kept distinct from any directed acyclic graph (DAG) extension used for Bayesian-network parameterization and inference. Logistic regression, a small neural network (MLP), and an RBF-kernel support-vector machine (SVM) provide discriminative baselines against the PC-learned Bayesian network. Classification is reported with AUROC, AUPRC, minority-class F1, and Brier score. Bayesian-network inference supports posterior queries from partial evidence, offering an explicit explanation path unavailable from a standard discriminative prediction alone. The dataset is small and observational; therefore, learned edges are treated as statistical dependencies rather than causal findings. In preliminary cross-validated results, logistic regression (AUROC 0.956) and the SVM (0.938) rank best, the PC-learned Bayesian network is intermediate (0.812, with AUPRC 0.966 but minority-class F1 of only 0.069), and the MLP is weakest (0.672). Reported metrics are preliminary, and structural sensitivity, orientation validation, and held-out Bayesian-network evaluation are identified as required validation steps.}%
{Probabilistic graphical models, PC algorithm, Bayesian networks, conditional independence, lung cancer prediction, support-vector machine, neural network, explainable AI.}
\label{ms2}

\section{Introduction}
Clinical and survey variables may be statistically related: symptoms can co-occur, behavioral factors can be associated, and several observations can jointly change the probability assigned to a target. A classifier can predict a label without representing these dependencies explicitly. A probabilistic graphical model (PGM) represents variables and their factorization in a graph, and can answer probability queries when only some variables are observed.

This project studies the Lung Cancer Survey Dataset distributed through Kaggle \cite{kaggle}, which is also the dataset used in the lung-cancer appendix of the Neural Graphical Models paper \cite{ngm}. The dataset has 309 rows and 16 columns: demographic and symptom/behavior variables, including \var{AGE}, and the binary target \var{LUNG\_CANCER}. A preliminary data audit found 33 exact duplicate rows, no missing cells, and 270 positive versus 39 negative target labels before deduplication. Duplicate removal yields 276 unique rows (238 YES, 38 NO). This strong class imbalance and small sample size materially constrain interpretation.

The objectives are:
\begin{enumerate}\itemsep0pt
\item Learn a conditional-independence skeleton using the PC algorithm and orient only edges justified by PC orientation rules.
\item Learn a Bayesian Network (BN) parameterization for probabilistic queries.
\item Compare BN predictions with logistic regression, a small neural network, and an RBF-kernel SVM using identical stratified evaluation folds.
\item Report classification performance using metrics sensitive to class imbalance and probability quality.
\item Demonstrate explainability through partial-evidence posterior inference and graph-based dependency inspection.
\end{enumerate}

\section{Problem Formulation}
Let $V=\{X_1,\ldots,X_{15},Y\}$ be the set of observed variables, where $Y$ is \var{LUNG\_CANCER}. The target takes values in $\{0,1\}$. The survey features are mostly binary; \var{AGE} is continuous in the raw data and is discretized into three bins for categorical graphical-model learning. The age cut points must be fitted on the training partition only in predictive evaluation. A Bayesian Network is a DAG $G=(V,E)$ with joint distribution
\begin{equation}
P(V)=\prod_{i=1}^{|V|}P\big(X_i\mid \mathrm{Pa}(X_i)\big),
\end{equation}
where $\mathrm{Pa}(X_i)$ denotes the parents of node $X_i$. For evidence $e$ and target value $y$, inference returns
\begin{equation}
P(Y=y\mid e)=\frac{P(Y=y,e)}{\sum_{y'}P(Y=y',e)}.
\end{equation}
The project is observational. An edge in the learned graph does not, by itself, establish a causal effect.

\section{Dataset and Preprocessing}
The principal dataset is the Kaggle Lung Cancer Survey Dataset \cite{kaggle}. The initial file has 309 records and 16 columns. The target distribution before deduplication is 270 YES and 39 NO; 33 exact duplicate rows were detected (Table~\ref{tab:m2audit}). Exact duplicates are removed before cross-validation to prevent identical rows appearing in both training and test partitions. The resulting dataset has 276 records (238 YES, 38 NO), so deduplication removed 32 YES rows and one NO row. \var{AGE} ranges from 21 to 87 (mean 62.7, median 62). The NGM paper \cite{ngm} describes the same Kaggle file as containing 284 patients with \var{AGE} binarized; the 309-row version used here differs in row count, so numbers are not directly comparable.

\begin{table}[!htbp]
\centering
\caption{Dataset audit before model fitting.}
\label{tab:m2audit}
\tablesize
\begin{tabular}{lc}
\toprule
Item & Value \\
\midrule
Raw rows & 309 \\
Columns & 16 \\
Exact duplicate rows removed & 33 \\
Rows after deduplication & 276 \\
Missing cells reported & 0 \\
Raw target YES & 270 \\
Raw target NO & 39 \\
Target YES / NO after deduplication & 238 / 38 \\
\bottomrule
\end{tabular}
\end{table}

Column names are normalized. \var{GENDER} is encoded as a binary variable; symptom/behavior codes are mapped consistently from the dataset's 1/2 convention to 0/1; the target is mapped to 0/1. The pipeline should fail loudly if unexpected values or missing cells appear. For predictive cross-validation, all learned transformations---including age-bin cut points, scaling, and probability calibration---must be fitted on training folds only.

\section{Background: PC Algorithm}
The PC algorithm learns a graph from conditional-independence (CI) tests rather than optimizing a graph score \cite{spirtes,colombo}. Its two phases are skeleton discovery and orientation.

\subsection{Skeleton Discovery}
The algorithm initializes a complete undirected graph. At conditioning-set size $\ell=0$ it tests marginal independence for every adjacent pair. It then increases $\ell$ and tests adjacent pairs conditioned on subsets of the current neighbors of one endpoint. If a test fails to reject conditional independence at significance level $\alpha$, the edge is removed and the conditioning set that produced the removal is stored as the separating set $\mathrm{Sep}(X,Y)$.

\subsection{Orientation}
For an unshielded triple $X-Z-Y$ where $X$ and $Y$ are not adjacent, if $Z\notin\mathrm{Sep}(X,Y)$ PC orients the triple as a collider: $X\to Z\leftarrow Y$. Meek orientation rules are then applied repeatedly to orient additional edges without introducing unsupported colliders or directed cycles. The standard output is a CPDAG, representing a Markov-equivalence class. Undirected edges are meaningful: the available observational evidence may not identify their directions. A DAG extension, if needed for BN inference, must be documented separately and checked for acyclicity.

\subsection{CI Tests and Finite-Sample Concerns}
For categorical variables, a common CI test is a conditional Pearson chi-square test. Sparse contingency tables can invalidate the usual asymptotic approximation and reduce test power. The small dataset and high-dimensional conditioning sets make this concern important. The significance level $\alpha$, maximum conditioning-set size, test count, and separating sets must be saved as experiment outputs. Sensitivity analysis across multiple $\alpha$ values is required before making strong structural claims.

\section{Learning Pipeline}
\subsection{PC Structural Learning}
The implementation should: (1) start from a complete undirected graph over all variables; (2) test candidate pairs in increasing conditioning-set size; (3) remove an edge only after a CI test supports conditional independence at the configured threshold; (4) store the exact separating set associated with every removed edge; (5) orient unshielded colliders using the stored separating sets; (6) apply Meek orientation rules to a fixed point, preventing directed cycles; and (7) export the skeleton, separating-set table, CPDAG, orientation log, and any DAG extension as separate artifacts. A visual graph alone is not sufficient evidence of correctness. Implementation checks should include symmetry of the skeleton, no self-edges, correct use of separating sets, and acyclicity of any fully directed extension.

\subsection{Learned Structure and the SMOKING Node}
Fig.~\ref{fig:m2skel} shows the undirected PC conditional-independence skeleton, and Fig.~\ref{fig:m2dag} shows the DAG extension used for Bayesian-network parameterization and inference. In both, the target \var{LUNG\_CANCER} is red and its Markov blanket / direct neighbours---\var{ALLERGY}, \var{ANXIETY}, \var{PEER\_PRESSURE}, \var{SWALLOWING\_DIFFICULTY} and \var{SHORTNESS\_OF\_BREATH}---are yellow. In the DAG, \var{ALLERGY} points into the target and the target points into \var{SWALLOWING\_DIFFICULTY}. \var{AGE} has no edges and is isolated in both graphs. Edges indicate statistical dependencies retained by the learning procedure, not causal directions.

\begin{figure}[!htbp]
\centering
\includegraphics[width=0.85\columnwidth]{figs/skeleton.png}
\caption{PC conditional-independence skeleton (undirected). \var{LUNG\_CANCER} is red; yellow nodes are its Markov blanket / direct neighbours. Thick red edges are the four \var{SMOKING} edges added by the clinical merge, not discovered by PC.}
\label{fig:m2skel}
\end{figure}

\emph{SMOKING merge.} In the PC run itself, every edge incident to \var{SMOKING} was removed during skeleton discovery. Because this contradicts clinical expectation, \var{SMOKING} was ``protected'' in the plotted graphs: four edges, to \var{WHEEZING}, \var{YELLOW\_FINGERS}, \var{CHRONIC\_DISEASE} and \var{COUGHING}, were added back by a clinical merge (thick red edges). These four edges are therefore a modeling choice, not a PC finding, and in the DAG extension all four are drawn pointing into \var{SMOKING}, an orientation that reflects the extension procedure rather than a causal claim. The recorded CI log (Table~\ref{tab:m2smoking}) lists removal decisions involving \var{SMOKING}. Under $\alpha=0.05$, each reported $p$-value exceeds the threshold, so the removals are consistent with the decision rule; the first three are close to the threshold and may be sensitive to sampling variation or test configuration.

\begin{table}[!htbp]
\centering
\caption{Recorded PC decisions involving \var{SMOKING}.}
\label{tab:m2smoking}
\tablesize
\setlength{\tabcolsep}{4pt}
\begin{tabular}{llc}
\toprule
Pair & Conditioning set & $p$-value \\
\midrule
\var{SMOKING}--\var{ANXIETY} & \var{COUGHING} & 0.06196 \\
\var{SMOKING}--\var{CHRONIC\_DISEASE} & \var{GENDER} & 0.05313 \\
\var{SMOKING}--\var{WHEEZING} & \var{FATIGUE} & 0.05288 \\
\var{SMOKING}--\var{COUGHING} & \var{ANXIETY} & 0.11486 \\
\bottomrule
\end{tabular}
\end{table}

The removal of the \var{SMOKING} edges must not be interpreted as evidence that smoking has no statistical or medical relationship with lung cancer. The PC algorithm's decisions depend on the sample, the CI test, the significance level, the conditioning-set search, variable encoding, and assumptions such as faithfulness. In a small dataset, sparse contingency tables can also make asymptotic chi-square $p$-values unreliable. Sensitivity analysis---repeating skeleton learning at $\alpha\in\{0.01,0.05,0.10\}$, inspecting expected cell counts, and comparing edge stability across resampled datasets---is the appropriate next step; until then the \var{SMOKING} result is a preliminary, configuration-dependent finding.

\subsection{PC-Learned Bayesian Network}
Fig.~\ref{fig:m2dag} shows the Bayesian network over all 16 variables obtained by extending the PC output to a DAG, with the \var{SMOKING} merge applied. This is the structure underlying the ``PC-learned Bayesian Network'' classifier in Section~\ref{sec:m2res}. The skeleton and the DAG extension are shown separately so that unresolved CPDAG orientations are not mistaken for learned causal directions.

\begin{figure}[!htbp]
\centering
\includegraphics[width=0.85\columnwidth]{figs/dag_m2.png}
\caption{PC-learned Bayesian network (DAG extension, \var{SMOKING} merged). \var{LUNG\_CANCER} is red, its Markov blanket / direct neighbours yellow. Arrows are the DAG orientation and are not causal directions.}
\label{fig:m2dag}
\end{figure}

\subsection{Bayesian-Network Parameters}
Given a DAG extension, the categorical conditional probability tables can be estimated with Dirichlet smoothing. For a node with $r_i$ states, parent configuration $j$ and state count $N_{ijk}$, the posterior-mean estimate is
\begin{equation}
\hat\theta_{ijk}=\frac{N_{ijk}+\alpha_0}{N_{ij}+r_i\alpha_0},
\end{equation}
where $\alpha_0>0$ is the per-state smoothing parameter. Smoothing avoids zero probabilities for unseen configurations. The parent-set size should be controlled because the sample size is small.

\subsection{Inference}
For evidence $e$, exact variable elimination computes the target posterior by multiplying local factors, summing out non-query, non-evidence variables, and normalizing the remaining factor. The implementation should be tested against brute-force enumeration on small configurations, with a numerical tolerance documented in the test output. Queries should include $P(\text{\var{LUNG\_CANCER}}=\text{YES}\mid e)$ for multiple partial-evidence sets; posterior changes when a single evidence variable is added or removed; the target's Markov blanket and its relation to the local factorization; and inference runtime as the number of observed variables changes. Any numerical posterior in the final results must be copied from an executed inference output, together with its evidence configuration, and must not be presented as a clinical risk estimate.

\section{Motivation for Neural Graphical Models}
A conventional BN represents conditional dependencies explicitly and supports probabilistic inference, but discrete conditional probability tables can become large and data-hungry when parent configurations proliferate. A conventional neural network can capture nonlinear predictive patterns but does not inherently provide a graphical factorization or general-purpose inference over missing variables. Neural Graphical Models (NGMs) seek to combine graph-structured inductive bias with neural parameterizations, enabling flexible nonlinear local models while retaining structure-aware learning and inference \cite{ngm}. This is a motivation for the broader research direction, not a claim that the present implementation reproduces the full NGM method. Unless a graph-constrained NGM architecture is implemented and evaluated, NGM results from the literature are reported only as an external reference.

\section{Classification Baselines}
\subsection{Neural Network}
A compact multilayer perceptron with regularization and early stopping or a fixed training protocol. Given the small dataset, a large architecture is avoided. Scaling and all preprocessing are fitted within each training fold, and random seeds and architecture settings are recorded.

\subsection{Support-Vector Machine}
An RBF-kernel SVM. Scaling is mandatory and is fitted inside each training fold. If probability estimates are enabled, probability calibration must be performed without using the held-out fold, and the exact hyperparameters and calibration method are reported.

\subsection{Bayesian Network Classifier}
For each held-out case, compute $P(Y=1\mid X_{\mathrm{observed}})$ using a BN whose structure and parameters were learned without the held-out case. If a single DAG extension is used, state how unresolved CPDAG edges were extended. A BN fitted to the full dataset may be used for illustrative inference examples, but those examples must not be presented as held-out predictive performance.

\section{Evaluation Protocol and KPIs}
Stratified cross-validation with identical fold assignments is used for all predictive models. With only 39 negative cases before deduplication, the number of folds must leave each test fold enough minority-class observations; fold-level metrics and mean $\pm$ standard deviation are reported, and any tuning must occur within training data. The four primary classification metrics are: (1) AUROC, threshold-independent ranking quality; (2) AUPRC, positive-class retrieval quality under class imbalance; (3) F1 for the NO class, with the positive label explicitly set to NO; and (4) Brier score, the mean squared error of predicted probabilities (lower is better). Confusion matrices and calibration plots are supporting diagnostics. For structure learning, node/edge counts, graph density, directed/undirected edge counts, number of CI tests, maximum conditioning-set size reached and sensitivity of edge inclusion to $\alpha$ are reported; for inference, numerical validation error and query latency.

\section{Preliminary Results and Interpretation}
\label{sec:m2res}
\subsection{Classifier Comparison}
The preliminary comparison is given in Table~\ref{tab:m2kpi} and Fig.~\ref{fig:m2kpi}, as mean scores across identical stratified folds. The values were transcribed from the saved comparison chart and are accurate to about $\pm0.005$; they are not final estimates until the exact protocol, fold-level scores, preprocessing and random seeds are archived and independently rerun.

\begin{table}[!htbp]
\centering
\caption{Preliminary classification metrics, mean over identical stratified folds (not final validated estimates).}
\label{tab:m2kpi}
\tablesize
\setlength{\tabcolsep}{5pt}
\begin{tabular}{lcccc}
\toprule
Model & AUROC & AUPRC & F1 (NO) & Brier \\
\midrule
Logistic Regression & 0.956 & 0.993 & 0.645 & 0.083 \\
Neural Network (MLP) & 0.672 & 0.923 & 0.137 & 0.173 \\
PC-learned BN & 0.812 & 0.966 & 0.069 & 0.102 \\
SVM (RBF) & 0.938 & 0.990 & 0.504 & 0.075 \\
\bottomrule
\end{tabular}
\end{table}

\begin{figure}[!htbp]
\centering
\includegraphics[width=\columnwidth]{figs/kpi_m2.png}
\caption{Classifier KPIs (AUPRC, AUROC, Brier score, F1 for the NO class), mean across identical stratified folds. Higher is better except for Brier score. Values are preliminary.}
\label{fig:m2kpi}
\end{figure}

Logistic regression and the RBF-SVM lead on ranking quality (AUROC 0.956 and 0.938), and the SVM has the lowest Brier score (0.075). The PC-learned BN is intermediate on AUROC (0.812) and keeps a high AUPRC (0.966), but its minority-class F1 is the lowest of all models (0.069), meaning it rarely identifies the NO class correctly at the default threshold. The MLP is weakest on AUROC (0.672) and Brier score (0.173) and also has poor minority-class F1 (0.137). High AUPRC for every model largely reflects the 238:38 class ratio; AUROC and F1 for NO are more informative here. The BN's gap to logistic regression is consistent with a sparse PC structure discarding predictive information (for example, \var{SMOKING} loses all PC edges), but this explanation has not been tested.

\subsection{SMOKING Node}
Each recorded $p$-value in Table~\ref{tab:m2smoking} exceeds $\alpha=0.05$ so the decisions are internally consistent with that threshold, but three are close to it; this is not robust evidence that smoking is independent of every variable. The graphs contain four \var{SMOKING} edges only because of the clinical merge. A sensitivity analysis at multiple thresholds and review of the CI-test implementation are required before interpreting the removed \var{SMOKING} edges substantively.

\subsection{Caveats}
These results describe the behavior of one run, not a definitive scientific conclusion. The dataset is self-reported, small, imbalanced, and not a clinical cohort suitable for establishing causal risk. The variables are not necessarily temporally ordered, and symptoms may be consequences or correlates of disease rather than pre-diagnostic predictors. Leakage and dataset-construction artifacts are additional concerns.

\section{Explainability Advantages of BN Inference}
A discriminative model typically maps a fixed feature vector to a class score. A BN instead represents a joint distribution and can answer conditional queries for different evidence subsets. The posterior $P(Y\mid e)$ explicitly identifies what evidence was provided and which variables were marginalized out. The graph and the target's Markov blanket help communicate local dependency structure. A useful report example presents the query, the evidence values, the computed posterior for both target states, the predicted class at a stated threshold, and a comparison with a changed evidence set. This is a model-based explanation of the model's probability calculation, not a causal explanation of disease, and the posterior must not be presented as medically validated individual risk.

\section{Motivation and Position Relative to SOTA}
The Neural Graphical Models paper by Shrivastava and Chajewska is the selected research anchor \cite{ngm}. Given data $X$ and a dependency graph $G$ (directed, undirected or mixed; expert-supplied or recovered by an algorithm such as uGLAD), an NGM fits a multilayer perceptron as a multi-task regression, so that each feature is a function of its neighbors (its Markov blanket for a DAG). Graph adherence is enforced through the path-norm product of the absolute weight matrices, giving the objective
\begin{equation}
\min_{\mathcal W,\mathcal B}\ \sum_{k=1}^{M}\big\|X^k-f_{\mathcal W,\mathcal B}(X^k)\big\|_2^2+\lambda\log\Big\|\Big(\prod_{i=1}^{L}|W_i|\Big)\ast S^c\Big\|_1,
\end{equation}
where $S^c$ is the complement of the graph adjacency matrix and $\ast$ is the Hadamard product \cite{ngm}. Inference uses message passing or gradient-based MAP updates of the unknown inputs, and sampling proceeds feature by feature along a BFS or topological order.

The paper's lung-cancer appendix uses the same Kaggle survey: a uGLAD conditional-independence graph, a 2-layer NGM with hidden size $H=30$ and tanh activation, and 5-fold cross-validation (Table~\ref{tab:m2ngm}). The scores are external results taken from the paper, whose table does not state the metric; they are not produced by this project.

\begin{table}[!htbp]
\centering
\caption{External reference: 5-fold CV results reported in \cite{ngm} for the lung-cancer data (metric not stated).}
\label{tab:m2ngm}
\tablesize
\begin{tabular}{lcc}
\toprule
Method & Lung cancer & Smoking \\
\midrule
Logistic regression & 0.95$\pm$0.02 & 0.71$\pm$0.01 \\
NGM & 0.96$\pm$0.01 & 0.79$\pm$0.02 \\
\bottomrule
\end{tabular}
\end{table}

The paper's NGM improves on logistic regression for predicting \var{SMOKING} (0.79 versus 0.71), the node whose edges PC removed here, which suggests that nonlinear local models may recover dependencies a sparse discrete BN misses. Our logistic regression on \var{LUNG\_CANCER} (AUROC 0.956) is numerically close to the paper's 0.95, but because the metric, row count (284 versus 309) and \var{AGE} handling differ, we draw no direct comparison. The current project does not implement an NGM; training a graph-constrained NGM on the PC graph of Fig.~\ref{fig:m2dag} is the natural next experiment.

\section{Limitations and Required Validation}
\begin{itemize}\itemsep0pt
\item \textbf{Small sample size:} 276 unique rows are insufficient for stable high-dimensional structure discovery without strong assumptions.
\item \textbf{Class imbalance:} the negative class is small; AUROC must be read alongside AUPRC, minority-class F1, and Brier score.
\item \textbf{CI-test validity:} asymptotic chi-square tests may be unreliable for sparse conditional tables; alternative tests or careful category/conditioning limits should be investigated.
\item \textbf{Structural sensitivity:} edge stability must be evaluated over significance levels and bootstrap/subsample runs.
\item \textbf{CPDAG versus DAG:} unresolved orientations must remain unresolved in the CPDAG; any DAG extension is an additional modeling choice.
\item \textbf{Predictive evaluation:} BN predictions must be computed on held-out folds using fold-specific structure and parameter learning before claiming a cross-validated BN score.
\item \textbf{Causality:} observational associations and learned directions do not establish causal effects without additional assumptions and evidence.
\item \textbf{External validity:} the survey dataset is not a representative clinical cohort; predictions are not suitable for medical diagnosis or screening.
\end{itemize}

\section{Conclusion}
This project combines PC-based conditional-independence structure learning, Bayesian-network probabilistic inference, and predictive baselines based on a neural network and an RBF SVM. The principal methodological requirement is to preserve the distinction between the PC skeleton, the oriented CPDAG, and any DAG extension used for inference. On the metrics that account for class imbalance and probability quality, logistic regression and the SVM outperformed the PC-learned BN and the MLP in the current preliminary run. BN inference provides a meaningful explainability advantage by supporting transparent posterior queries under partial evidence, but it does not establish causal relationships. The present numbers are preliminary; final conclusions require reproducible cross-validation, structural sensitivity analysis, validation of PC orientation rules, and held-out BN inference evaluation.

\section*{Reproducibility Checklist}
\noindent Before submission, archive the exact source revision, dependency versions, random seeds, dataset checksum, PC significance level, CI-test choice, conditioning-set limit, all separating sets, CPDAG and DAG outputs, cross-validation fold assignments, fold-level predictions, and scripts that generate each table and figure.

\msheader{Milestone 3: Results and Analysis}%
{This report presents the Milestone 3 (M3) results of our study of lung-cancer risk factors with probabilistic graphical models (PGMs). On 276 de-duplicated survey records we ran $5\times5$ repeated stratified cross-validation and compared a PC-learned Bayesian network (BN), with and without four clinically protected \var{SMOKING} edges, against logistic regression, an RBF-kernel SVM and a small MLP under identical folds. Logistic regression discriminated best (AUROC $0.938\pm0.009$), followed by the SVM ($0.912\pm0.020$), the BN ($0.822\pm0.008$) and the MLP ($0.566\pm0.127$). The BN had the lowest expected calibration error (0.043) but no better Brier score or log-loss than logistic regression or the SVM, and at the 0.5 threshold it recognised only 12.6\% of NO cases (F1 for NO $0.187\pm0.080$). The protected edges changed nothing measurable. For structure learning, the PC skeleton grew from 7.3 to 18 edges as the training fraction rose from 20\% to 100\% without plateauing; some edges incident to the target were unstable. Exact inference matched brute-force enumeration to $2.2\times10^{-16}$ and ran in 0.10--0.42\,ms. We conclude that on this dataset the PGM's advantage is explicit partial-evidence probabilistic inference and an inspectable dependency graph, not discriminative accuracy. No NGM was implemented, so literature numbers remain external references.}%
{Probabilistic graphical models, PC algorithm, Bayesian networks, calibration, bootstrap edge stability, lung cancer, robustness, statistical comparison.}
\label{ms3}

\section{Introduction}
Milestone~1 asked whether a PGM that represents dependencies among risk factors and symptoms can match the predictive quality of standard classifiers while adding an interpretable dependency map and calibrated uncertainty. Milestone~2 built the pipeline: data audit, PC structure learning, a Dirichlet-smoothed BN with exact inference, and three discriminative baselines. M3 asks what the evidence says. Following the course brief, our aim is not to make the PGM win but to determine when it helps and when it does not.

We test three hypotheses stated or implied in M1:
\begin{itemize}\itemsep0pt
\item \textbf{H1 (prediction).} A PC-learned BN predicts \var{LUNG\_CANCER} about as well as logistic regression, an SVM or an MLP.
\item \textbf{H2 (structure).} PC recovers a stable, interpretable dependency structure from 276 samples.
\item \textbf{H3 (inference).} The BN provides exact, fast, calibrated posteriors under partial evidence, which fixed-input classifiers cannot.
\end{itemize}
This report adds over M2: repeated cross-validation with mean$\pm$SD and 95\% intervals; paired tests; log-loss, Brier score, expected calibration error (ECE) and reliability diagrams; a protected-edge ablation; $\alpha$-sensitivity, bootstrap edge stability and sample-size studies of the learned structure; a graphical-lasso comparison; test-time missingness and noise; exact-inference validation and runtime; error analysis; and a revised SOTA position relative to Neural Graphical Models (NGMs) \cite{ngm}. All numbers are taken from the saved experiment outputs; any quantity computed afterwards from those files is stated as such.

\section{Experimental Design}
\subsection{Data and Preprocessing}
The Kaggle lung-cancer survey \cite{kaggle} has 309 rows and 16 columns; 33 exact duplicates were removed, leaving 276 rows (238 YES, 38 NO). The majority-class rate is 0.862, so a trivial always-YES classifier has accuracy 0.862. \var{AGE} spans 21--87. For the BN, \var{AGE} is discretised into three bins whose cut points are fitted on the training fold only; the other models use a median-imputed, standardised \var{AGE}. All learned transformations are fitted inside training folds.

\subsection{Models and Settings}
Logistic regression (LR): balanced class weights, \texttt{max\_iter}=3000. SVM: RBF kernel, $C=1$, balanced class weights, Platt-scaled probabilities. MLP: hidden layers (32,16), L2 penalty 0.01, early stopping, \texttt{max\_iter}=1200. PC-learned BN: skeleton by PC with a conditional likelihood-ratio ($G$) test summed over conditioning strata, $\alpha=0.05$, maximum conditioning-set size 3; collider and Meek orientation; DAG extension; Dirichlet-smoothed CPTs ($\alpha_0=1$); exact variable elimination. Structure and parameters are re-learned from each training fold. Two BN variants are compared: plain PC and protected edges, in which four \var{SMOKING} edges (to \var{WHEEZING}, \var{YELLOW\_FINGERS}, \var{CHRONIC\_DISEASE}, \var{COUGHING}) are forced into the skeleton as a clinical modelling choice, as in M2. No hyperparameters were tuned for any model.

\subsection{Protocol and Statistics}
We use RepeatedStratifiedKFold (5 folds $\times$ 5 repeats, seed 516) with identical splits for all models. Out-of-fold predictions are pooled within each repeat (276 predictions), giving five repeat-level observations per model, from which we report mean $\pm$ SD (ddof=1) and a 95\% $t$-interval (df=4). F1 for the NO class uses the 0.5 threshold with NO as the positive label. ECE uses ten equal-width bins. Paired comparisons use a paired $t$-test and the Wilcoxon signed-rank test across the five repeats. These are exploratory: repeats reuse the same 276 rows, so they are not independent and $t$-test $p$-values are optimistic \cite{nadeau}. Moreover, with $n=5$ pairs the smallest attainable two-sided Wilcoxon $p$ is 0.0625, so it cannot reach 0.05 by construction; we therefore rely on effect sizes and consistency across repeats (wins out of 5) as the primary evidence.

\subsection{Further Experiments}
(i) PC at $\alpha\in\{0.01,0.05,0.10\}$, plain and protected; (ii) 100 bootstrap resamples at $\alpha=0.05$ with plain PC; (iii) subsamples of 20--100\% of the data (10 draws each); (iv) test-time corruption of held-out features (10\%/20\% MCAR; 5\%/10\% feature noise, i.e.\ binary flips and Gaussian noise on \var{AGE}) with models trained on clean data and train-fold imputation; (v) graphical lasso with cross-validated penalty; (vi) variable elimination (VE) against brute-force enumeration with 30 random queries per evidence-set size; (vii) pooled confusion matrices. No structural Hamming distance (SHD) is reported because the real data have no ground-truth graph.

\section{Predictive Results}
\subsection{KPI Comparison}
Table~\ref{tab:m3kpi} and Fig.~\ref{fig:m3kpi} give the main results; Table~\ref{tab:m3ci} gives 95\% confidence intervals.

\begin{table}[!htbp]
\centering
\caption{Repeated-CV KPIs (mean $\pm$ SD over 5 repeats of 5-fold CV; best per column in bold).}
\label{tab:m3kpi}
\tablesize
\setlength{\tabcolsep}{4pt}
\resizebox{\columnwidth}{!}{%
\begin{tabular}{lccccccc}
\toprule
Model & AUROC $\uparrow$ & AUPRC $\uparrow$ & F1 (NO) $\uparrow$ & Brier $\downarrow$ & Log-loss $\downarrow$ & ECE $\downarrow$ & Accuracy \\
\midrule
Logistic Regression & \textbf{0.938$\pm$0.009} & \textbf{0.990$\pm$0.002} & \textbf{0.623$\pm$0.027} & 0.091$\pm$0.005 & 0.301$\pm$0.013 & 0.118$\pm$0.003 & 0.864$\pm$0.008 \\
SVM (RBF) & 0.912$\pm$0.020 & 0.984$\pm$0.005 & 0.508$\pm$0.072 & \textbf{0.079$\pm$0.005} & \textbf{0.253$\pm$0.016} & 0.052$\pm$0.011 & \textbf{0.882$\pm$0.011} \\
PC-BN (plain PC) & 0.822$\pm$0.008 & 0.967$\pm$0.003 & 0.187$\pm$0.080 & 0.099$\pm$0.004 & 0.316$\pm$0.010 & 0.043$\pm$0.005 & 0.854$\pm$0.015 \\
PC-BN (protected edges) & 0.820$\pm$0.009 & 0.967$\pm$0.003 & 0.170$\pm$0.080 & 0.100$\pm$0.003 & 0.317$\pm$0.010 & \textbf{0.040$\pm$0.006} & 0.854$\pm$0.008 \\
Neural Network (MLP) & 0.566$\pm$0.127 & 0.883$\pm$0.051 & 0.085$\pm$0.121 & 0.191$\pm$0.021 & 0.572$\pm$0.045 & 0.268$\pm$0.034 & 0.843$\pm$0.015 \\
\bottomrule
\end{tabular}}
\\[2pt]{\scriptsize Majority-class (always YES) accuracy: 0.862.}
\end{table}

\begin{figure}[!htbp]
\centering
\includegraphics[width=\columnwidth]{figs/kpigrid.png}
\caption{Classifier KPIs, mean $\pm$ SD across 5 repeats of 5-fold CV (identical splits for all models). BN = PC-learned Bayesian network; ``prot.'' = with protected \var{SMOKING} edges. Values are those in Table~\ref{tab:m3kpi}.}
\label{fig:m3kpi}
\end{figure}

\begin{table}[!htbp]
\centering
\caption{95\% confidence intervals ($t$, df=4) for the repeat-level means.}
\label{tab:m3ci}
\tablesize
\setlength{\tabcolsep}{2.5pt}
\begin{tabular}{lccccc}
\toprule
Model & AUROC & F1 (NO) & Brier & Log-loss & ECE \\
\midrule
LR & [.926,.949] & [.590,.657] & [.085,.097] & [.285,.318] & [.114,.121] \\
SVM & [.887,.936] & [.418,.597] & [.073,.085] & [.233,.273] & [.038,.067] \\
BN plain & [.812,.832] & [.087,.287] & [.094,.105] & [.304,.328] & [.036,.049] \\
BN prot. & [.809,.832] & [.071,.269] & [.097,.104] & [.305,.329] & [.033,.047] \\
MLP & [.408,.724] & [.000,.234]$^*$ & [.165,.217] & [.516,.628] & [.225,.310] \\
\bottomrule
\end{tabular}
\\[2pt]{\scriptsize $^*$Lower limit $-0.065$ truncated at 0.}
\end{table}

\emph{What happened.} LR has the highest AUROC, AUPRC and F1(NO); the SVM has the lowest Brier score and log-loss and the highest accuracy; the BN is third on AUROC; the MLP is last on almost every metric. The AUROC intervals of LR and the SVM overlap slightly, whereas the BN interval ([0.812, 0.832]) is disjoint from both. The MLP is also highly unstable (AUROC SD 0.127; individual repeats ranged from 0.383 to 0.667).

\emph{Was it expected?} The ordering LR $>$ SVM $>$ BN $>$ MLP agrees with M2, with all values moving slightly down (e.g.\ LR 0.956$\to$0.938, MLP 0.672$\to$0.566), which is what repeating the split should do after a single favourable partition. AUPRC is high for every model (0.88--0.99) mainly because 86\% of rows are YES; it is not informative here. Likewise accuracy is nearly uninformative: LR (0.864) equals the always-YES rate, the BN (0.854) and MLP (0.843) are below it, and only the SVM (0.882) is above it.

\emph{Why.} In the final DAG, the Markov blanket of \var{LUNG\_CANCER} contains only five variables (\var{ALLERGY}, \var{PEER\_PRESSURE}, \var{SWALLOWING\_DIFFICULTY}, \var{ANXIETY}, \var{SHORTNESS\_OF\_BREATH}), so the BN classifier effectively uses 5 of 15 features, whereas LR and the SVM use all. At the full-data CI test, several variables are strongly associated with the target marginally (e.g.\ \var{ALCOHOL\_CONSUMING}, $p=5.8\times10^{-7}$; \var{COUGHING}, $2.5\times10^{-5}$; \var{WHEEZING}, $2.8\times10^{-5}$) but are not in its blanket. This is consistent with a sparse PC structure discarding predictive information, but we did not run the experiment that would test it (e.g.\ a BN with the blanket expanded). The MLP's instability is probably an untuned-small-data effect (early stopping uses a validation split of only a few NO cases), but we did not diagnose it; we therefore treat it as a weak baseline rather than as evidence about neural models in general.

\emph{Limitation.} Pooling 276 out-of-fold predictions per repeat gives stable estimates of the mean but only five repeat-level points; and only 38 NO cases limit every NO-class metric.

\subsection{Statistical Comparison}
Table~\ref{tab:m3paired} compares every model with the plain-PC BN across the five repeats.

\begin{table}[!htbp]
\centering
\caption{Paired comparison against the plain-PC BN over 5 repeats: difference of means (model $-$ BN); paired $t$-test $p$ / Wilcoxon $p$.}
\label{tab:m3paired}
\tablesize
\setlength{\tabcolsep}{6pt}
\resizebox{\columnwidth}{!}{%
\begin{tabular}{lcccc}
\toprule
Metric & LR & SVM & MLP & BN (protected) \\
\midrule
AUROC $\uparrow$ & $+0.115$; $<$.001 / .063 & $+0.090$; .001 / .063 & $-0.256$; .011 / .063 & $-0.002$; .554 / .625 \\
F1 (NO) $\uparrow$ & $+0.437$; $<$.001 / .063 & $+0.321$; $<$.001 / .063 & $-0.102$; .279 / .313 & $-0.016$; .204 / .500 \\
Brier $\downarrow$ & $-0.008$; .106 / .188 & $-0.020$; .003 / .063 & $+0.092$; $<$.001 / .063 & $+0.001$; .421 / .875 \\
Log-loss $\downarrow$ & $-0.014$; .188 / .313 & $-0.063$; .001 / .063 & $+0.256$; $<$.001 / .063 & $+0.001$; .487 / .875 \\
ECE $\downarrow$ & $+0.075$; $<$.001 / .063 & $+0.010$; .218 / .313 & $+0.225$; $<$.001 / .063 & $-0.003$; .468 / .625 \\
\bottomrule
\end{tabular}}
\end{table}

LR and the SVM beat the BN on AUROC and F1(NO) in 5 of 5 repeats (effects of $+0.09$ to $+0.12$ AUROC and $+0.32$ to $+0.44$ F1), and the MLP loses to the BN on AUROC in 5 of 5 repeats. The difference in Brier score between LR and the BN is small ($-0.008$) and not distinguishable from zero ($p_t=0.106$); the SVM's Brier and log-loss advantages over the BN ($-0.020$, $-0.063$) are consistent across repeats. Between LR and the SVM, LR is better on AUROC ($+0.026$, $p_t=0.013$) and F1(NO) ($+0.116$, $p_t=0.010$), the SVM on Brier ($-0.012$, $p_t=0.006$) and log-loss ($-0.049$, $p_t=0.002$); the Wilcoxon $p$ is 0.0625 for all four, the minimum possible. Because of the shared data and the attainable-$p$ floor, these are consistent effects across repeats rather than confirmed significance.

\subsection{Probability Quality and Calibration}
The task requires probabilities, not only labels. On ECE, the BN (0.040--0.043) and the SVM (0.052) are best, LR is worse (0.118) and the MLP is worst (0.268). On Brier score and log-loss, however, the SVM is best, LR second and the BN third, so the BN's calibration advantage is not matched by sharper or more accurate probabilities: its Brier (0.099) is within 0.008 of LR's and its log-loss (0.316) is 0.015 above it \cite{guo}.

Fig.~\ref{fig:m3rel} shows how the models differ. LR lies mostly above the diagonal, i.e.\ its positive probabilities are too low; this is plausible given balanced class weights, which shift probabilities toward the minority class (not tested separately). The SVM and BN fall below the diagonal mainly in the 0.65--0.9 range, i.e.\ they are over-confident for YES there, and lie above it at lower probabilities; with few cases per bin these curves are noisy. The BN has no points below a predicted probability of about 0.3: it essentially never assigns a low probability of cancer, which is why its F1(NO) is poor. The MLP curve is non-monotone and collapses in its highest bin. ECE with ten bins is dominated by the highly populated high-probability bins, where positives are common, so a low ECE here does not imply good discrimination. Hence H3's ``calibrated'' claim is supported only in the narrow ECE sense.

\begin{figure}[!htbp]
\centering
\includegraphics[width=0.85\columnwidth]{figs/reliab.png}
\caption{Reliability diagram of pooled out-of-fold predictions (10 bins). The dashed line is perfect calibration. The bins at the extremes contain few NO cases, so those points are noisy.}
\label{fig:m3rel}
\end{figure}

\subsection{Ablation: Clinically Protected SMOKING Edges}
The protected and plain BN variants are statistically indistinguishable on every metric (Table~\ref{tab:m3paired}, last column; AUROC 0.820 vs.\ 0.822, F1(NO) 0.170 vs.\ 0.187, all $p_t>0.2$). The protected-edge variant had a nominally lower ECE (0.040 vs.\ 0.043), but the difference is a third of one SD. Imposing the clinical prior neither helps nor hurts prediction; its value is purely representational.

\subsection{Error Analysis}
Table~\ref{tab:m3conf} pools all out-of-fold predictions (each row appears once per repeat, so counts are five times the underlying 38 NO and 238 YES) at the 0.5 threshold.

\begin{table}[!htbp]
\centering
\caption{Pooled confusion counts at threshold 0.5 (NO = class 0; $5\times$ repeated rows).}
\label{tab:m3conf}
\tablesize
\setlength{\tabcolsep}{4pt}
\begin{tabular}{lcccccc}
\toprule
Model & TN & FP & FN & TP & Spec.\,(NO) & Sens.\,(YES) \\
\midrule
LR & 155 & 35 & 152 & 1038 & 0.816 & 0.872 \\
SVM & 85 & 105 & 58 & 1132 & 0.447 & 0.951 \\
BN plain & 24 & 166 & 35 & 1155 & 0.126 & 0.971 \\
BN prot. & 22 & 168 & 34 & 1156 & 0.116 & 0.971 \\
MLP & 11 & 179 & 38 & 1152 & 0.058 & 0.968 \\
\bottomrule
\end{tabular}
\end{table}

The BN, SVM and MLP behave like near-majority predictors: they find 12.6\%, 44.7\% and 5.8\% of NO cases while missing only 2.9\%, 4.9\% and 3.2\% of YES cases. LR makes a different trade, recognising 81.6\% of NO cases at the cost of missing 12.8\% of YES cases. Which error matters more is a clinical question; for screening, missed cancers are usually costlier, which would favour the BN/SVM operating point, but the dataset is not suitable for such a decision. The threshold was not tuned, so F1(NO) reflects the default operating point; the BN's AUROC of 0.82 shows it does have some ranking ability that a different threshold could expose. This was not evaluated.

\section{Structure-Learning Results}
\subsection{Learned Graph}
Fig.~\ref{fig:m3dag} shows the DAG learned on the full data with the protected edges. The skeleton has 22 edges (18 from plain PC plus the four \var{SMOKING} edges); the CPDAG has no undirected edges and no orientation conflicts, and \var{AGE} is isolated. The orientation log labels all 22 edges ``PC compelled''; for the four \var{SMOKING} edges this means only that they were oriented by the collider and Meek rules applied to the merged skeleton, not that PC discovered them. Because every edge ended up directed from a sparse graph with few separating-set decisions, orientation is a fragile output; we did not test orientation stability, so the arrows are not causal claims. The predictive BNs are re-learned per fold and can differ from this full-data graph.

\begin{figure}[!htbp]
\centering
\includegraphics[width=\columnwidth]{figs/dag.png}
\caption{PC-learned DAG extension on the full data ($\alpha=0.05$, protected \var{SMOKING} edges). \var{LUNG\_CANCER} is red; its Markov blanket is yellow; thick red edges are the four \var{SMOKING} edges added by the clinical merge. Arrows are DAG orientations, not causal directions.}
\label{fig:m3dag}
\end{figure}

\subsection{Sensitivity to the Significance Level}
Edge counts grow monotonically with $\alpha$ (Table~\ref{tab:m3alpha}, Fig.~\ref{fig:m3alpha}), and the edge sets are nested: every edge present at a smaller $\alpha$ is also present at the larger ones, for both variants. The neighbours of \var{LUNG\_CANCER} are unchanged between $\alpha=0.01$ and 0.05 and gain two edges at 0.10. The full-data run logged 852 CI tests (120, 402, 250 and 80 at conditioning-set sizes 0 to 3), so the size limit 3 was reached; plain PC at $\alpha=0.05$ used 761. Protected edges are still tested (75 logged tests) but never removed, so they and their neighbours stay in play longer. This matches the expectation that more permissive thresholds retain more weak dependencies, and shows that the target's local structure is stable across the usual $\alpha$ range while the global graph is not.

\begin{table}[!htbp]
\centering
\caption{Skeleton size versus $\alpha$ (density = edges/120 possible).}
\label{tab:m3alpha}
\tablesize
\setlength{\tabcolsep}{3pt}
\begin{tabular}{ccccc}
\toprule
$\alpha$ & Plain & Dens. & Prot. & Dens. \\
\midrule
0.01 & 14 & 0.117 & 18 & 0.150 \\
0.05 & 18 & 0.150 & 22 & 0.183 \\
0.10 & 25 & 0.208 & 26 & 0.217 \\
\bottomrule
\end{tabular}
\\[2pt]{\scriptsize Plain-PC target neighbours: \var{ALLERGY}, \var{PEER\_PRESSURE}, \var{SWALLOWING\_DIFFICULTY} at 0.01 and 0.05; plus \var{ALCOHOL\_CONSUMING}, \var{FATIGUE} at 0.10.}
\end{table}

\begin{figure}[!htbp]
\centering
\includegraphics[width=0.49\columnwidth]{figs/alpha.png}\hfill
\includegraphics[width=0.49\columnwidth]{figs/samplesize.png}
\caption{Left: skeleton edge count versus PC significance level $\alpha$ (plain vs.\ protected). Right: PC skeleton edges versus training-sample fraction (mean $\pm$ SD over 10 random subsamples, $\alpha=0.05$, plain PC).}
\label{fig:m3alpha}
\end{figure}

\subsection{Bootstrap Edge Stability}
Fig.~\ref{fig:m3boot} shows the 15 most frequently selected edges in 100 bootstrap resamples (plain PC, $\alpha=0.05$). Reading from the figure, eight edges are selected in roughly 90\% or more of resamples: \var{FATIGUE}--\var{SHORTNESS\_OF\_BREATH}, \var{ANXIETY}--\var{YELLOW\_FINGERS}, \var{ALCOHOL\_CONSUMING}--\var{GENDER}, \var{ANXIETY}--\var{SWALLOWING\_DIFFICULTY} (all about 1.0), \var{COUGHING}--\var{WHEEZING}, \var{CHEST\_PAIN}--\var{GENDER}, \var{ALLERGY}--\var{LUNG\_CANCER} and \var{COUGHING}--\var{SHORTNESS\_OF\_BREATH}. The remaining seven plotted edges lie roughly between 0.6 and 0.9. Among edges incident to the target, \var{ALLERGY}--\var{LUNG\_CANCER} is stable, \var{LUNG\_CANCER}--\var{SWALLOWING\_DIFFICULTY} is selected in about three quarters of resamples, and \var{PEER\_PRESSURE}--\var{LUNG\_CANCER}, although present in the full-data graph, does not appear in the top 15, so it was selected in fewer than roughly 60\% of resamples. The conclusion for H2 is therefore mixed: a stable core exists, but at least one of the target's three neighbours is not reliably recovered. We attach no clinical interpretation to the strongest edge (\var{ANXIETY}--\var{YELLOW\_FINGERS}); it is a statistical regularity of this survey. All frequencies are read from the figure.

\begin{figure}[!htbp]
\centering
\includegraphics[width=0.85\columnwidth]{figs/bootstrap.png}
\caption{Bootstrap selection frequency of the 15 most stable PC edges (100 resamples, plain PC, $\alpha=0.05$).}
\label{fig:m3boot}
\end{figure}

\subsection{Sample Size versus Learned Structure}
The number of learned edges rises steadily with $n$ (Table~\ref{tab:m3size}, Fig.~\ref{fig:m3alpha}) from 7.3 at 55 samples to 18 at 276, with no visible plateau between 80\% and 100\%. This is expected: CI tests gain power as $n$ grows, so weaker dependencies become detectable. It also implies that the 18-edge graph is itself sample-limited and that more data would probably add edges. A ``sample size versus SHD'' plot cannot be produced on real data without a ground-truth graph; edge count and density against $n$ are used instead. The 100\% point is a single deterministic run, and subsamples of smaller size are drawn from the same 276 rows, so the SDs understate sampling variability.

\begin{table}[!htbp]
\centering
\caption{PC skeleton versus sample size (plain PC, $\alpha=0.05$; 10 subsamples per size).}
\label{tab:m3size}
\tablesize
\setlength{\tabcolsep}{4pt}
\begin{tabular}{lccc}
\toprule
$n$ & Edges (mean$\pm$SD) & Density & CI tests \\
\midrule
55 (20\%) & 7.3$\pm$1.4 & 0.061 & 198$\pm$19 \\
110 (40\%) & 11.7$\pm$2.2 & 0.097 & 297$\pm$49 \\
166 (60\%) & 14.1$\pm$1.1 & 0.117 & 414$\pm$63 \\
221 (80\%) & 16.8$\pm$1.4 & 0.140 & 600$\pm$64 \\
276 (100\%) & 18 (single run) & 0.150 & 761 \\
\bottomrule
\end{tabular}
\end{table}

\subsection{Comparison with the Graphical Lasso}
The graphical lasso \cite{glasso} (cross-validated penalty 0.0276) returns 82 edges among 16 nodes (density 0.683) against 18--22 for PC (density 0.15--0.18). All 22 edges of the protected PC skeleton, and all 18 of plain PC, are among the 82 lasso edges, so PC's graph is a subgraph of the lasso graph (Jaccard similarity with the protected skeleton 0.27). Of the lasso edges, 37 have $|\rho|>0.1$ and 11 have $|\rho|>0.2$; the strongest are \var{YELLOW\_FINGERS}--\var{ANXIETY} (0.42), \var{FATIGUE}--\var{SHORTNESS\_OF\_BREATH} (0.33), \var{ANXIETY}--\var{SWALLOWING\_DIFFICULTY} (0.33), \var{GENDER}--\var{ALCOHOL\_CONSUMING} (0.29), \var{GENDER}--\var{CHEST\_PAIN} (0.24) and \var{WHEEZING}--\var{COUGHING} (0.24), which are also among PC's most stable edges. The two methods therefore agree on the strong core and disagree on the weak edges. They also disagree on the target: the lasso's largest partial correlations with \var{LUNG\_CANCER} are with \var{ALCOHOL\_CONSUMING} (0.21), \var{ALLERGY} (0.21) and \var{FATIGUE} (0.17), and PC links the target to \var{ALLERGY}, \var{SWALLOWING\_DIFFICULTY} and \var{PEER\_PRESSURE} instead. The lasso is a Gaussian-model estimate applied to mostly binary data and is not a causal graph; the two graphs are not directly comparable by a single score.

\subsection{The SMOKING Node}
In the full-data run PC removed every edge incident to \var{SMOKING}, each with $p>\alpha$ at the deciding conditioning set (Table~\ref{tab:m3smoking}). Three of the four recorded decisions lie within 0.013 of the threshold, and \var{SMOKING}--\var{LUNG\_CANCER} is marginally non-significant ($p=0.563$) in this sample. At $\alpha=0.10$ plain PC attaches \var{SMOKING} to \var{CHRONIC\_DISEASE}, \var{WHEEZING} and \var{ANXIETY}. The lasso keeps \var{SMOKING} connected to all variables but with small partial correlations (largest $|\rho|=0.136$ with \var{CHRONIC\_DISEASE}; 0.045 with \var{LUNG\_CANCER}). We read this as weak, threshold-dependent evidence in a small self-reported sample, not as evidence that smoking is unrelated to lung cancer. One caution: \var{SMOKING}--\var{YELLOW\_FINGERS} is marginally independent here ($p=0.73$), so protecting that edge imposes a prior that the data does not support. The $p$-values differ in the third decimal from the M2 table because the M3 run was re-executed with the final code.

\begin{table}[!htbp]
\centering
\caption{CI decisions involving \var{SMOKING} in the full-data run ($\alpha=0.05$).}
\label{tab:m3smoking}
\tablesize
\setlength{\tabcolsep}{3pt}
\begin{tabular}{llc}
\toprule
Pair & Conditioning set & $p$-value \\
\midrule
\var{SMOKING}--\var{ANXIETY} & \var{COUGHING} & 0.0630 \\
\var{SMOKING}--\var{CHRONIC\_DISEASE}$^\dagger$ & \var{GENDER} & 0.0523 \\
\var{SMOKING}--\var{WHEEZING}$^\dagger$ & \var{FATIGUE} & 0.0518 \\
\var{SMOKING}--\var{COUGHING}$^\dagger$ & \var{ANXIETY} & 0.1143 \\
\var{SMOKING}--\var{YELLOW\_FINGERS}$^\dagger$ & (marginal) & 0.7296 \\
\var{SMOKING}--\var{LUNG\_CANCER} & (marginal) & 0.5629 \\
\bottomrule
\end{tabular}
\\[2pt]{\scriptsize $^\dagger$Edge restored by the clinical merge.}
\end{table}

\section{Robustness Analysis}
Models were trained on clean data and the held-out features were corrupted (Table~\ref{tab:m3rob}, Fig.~\ref{fig:m3rob}); this uses a single 5-fold split, so each fold has about 55 test rows and about 8 NO cases. Missing values are imputed with training-fold medians/modes for all models, and the BN is the protected-edge variant. Fold-level metrics are averaged here, whereas Section~III pools predictions per repeat, which is why clean-condition values differ slightly (e.g.\ LR AUROC 0.958 here vs.\ 0.938; fold-level ECE is probably also inflated by the small test folds).

\begin{table}[!htbp]
\centering
\caption{Test-time corruption: AUROC (mean $\pm$ SD over 5 folds) and Brier score (mean).}
\label{tab:m3rob}
\tablesize
\setlength{\tabcolsep}{6pt}
\resizebox{\columnwidth}{!}{%
\begin{tabular}{lcccccccc}
\toprule
 & \multicolumn{4}{c}{AUROC} & \multicolumn{4}{c}{Brier (lower is better)} \\
\cmidrule(lr){2-5}\cmidrule(lr){6-9}
Condition & LR & SVM & BN & MLP & LR & SVM & BN & MLP \\
\midrule
Clean & .958$\pm$.023 & .939$\pm$.022 & .812$\pm$.041 & .553$\pm$.210 & .084 & .075 & .103 & .191 \\
MCAR 10\% & .948$\pm$.029 & .933$\pm$.030 & .786$\pm$.059 & .536$\pm$.224 & .079 & .080 & .110 & .192 \\
MCAR 20\% & .879$\pm$.103 & .869$\pm$.096 & .724$\pm$.069 & .637$\pm$.128 & .089 & .098 & .120 & .190 \\
Noise 5\% & .937$\pm$.021 & .923$\pm$.020 & .758$\pm$.058 & .535$\pm$.235 & .098 & .081 & .111 & .192 \\
Noise 10\% & .884$\pm$.076 & .882$\pm$.053 & .730$\pm$.063 & .481$\pm$.124 & .125 & .094 & .115 & .197 \\
\bottomrule
\end{tabular}}
\end{table}

\begin{figure}[!htbp]
\centering
\includegraphics[width=\columnwidth]{figs/robust.png}
\caption{AUROC under test-time missingness (MCAR) and noise, mean over 5 folds.}
\label{fig:m3rob}
\end{figure}

\emph{What happened.} All models degrade as the corruption rate increases. From clean to 20\% MCAR, AUROC drops by 0.078 (LR), 0.070 (SVM) and 0.088 (BN); from clean to 10\% noise by 0.074, 0.056 and 0.082. The ranking LR $\gtrsim$ SVM $>$ BN $>$ MLP is preserved in every condition (LR and SVM are tied at 10\% noise). LR's Brier score deteriorates most under noise (0.084$\to$0.125), overtaking the BN's (0.103$\to$0.115) at 10\% noise. The MLP's apparent gain at 20\% MCAR (0.553$\to$0.637) is noise around chance.

\emph{Interpretation.} Fold-level SDs (up to 0.10) are larger than the differences in degradation between models, so we cannot claim that any model is more robust than another, and we do not claim that the BN's explicit joint model protects it against missing evidence in this experiment. A fair test of that advantage would marginalise missing variables in the BN rather than impute them, which we did not do; this is the most natural follow-up to H3. What the experiment does show is that discriminative accuracy is retained to within about 0.1 AUROC under 20\% missingness, and that the ordering is not an artefact of clean data.

\section{Inference Validation and Computational Cost}
For each evidence-set size from 0 to 15 observed variables we ran 30 random posterior queries on the full-data BN. For 0--4 evidence variables (150 queries) the VE posterior agreed with brute-force enumeration to a maximum absolute error of $2.2\times10^{-16}$, i.e.\ machine precision; beyond four variables the brute-force check was not run. Runtime (Fig.~\ref{fig:m3run}) falls from $0.418\pm0.023$\,ms with no evidence to $0.102\pm0.007$\,ms with all 15 variables observed. The decline is expected: each observed variable is clamped and removed from elimination, so factors shrink. This is a property of one 16-node network on one machine, not a scalability result, and we did not time the discriminative models, so we make no speed comparison. Together with Section~III-A, it fits an accuracy--efficiency picture: BN inference is exact and effectively free, but is the less accurate predictor. In the BN, posterior queries with partial evidence are exact and need no retraining, which is the qualitative advantage argued for in M1 (H3) and the main reason to keep a PGM in the comparison.

\begin{figure}[!htbp]
\centering
\includegraphics[width=\columnwidth]{figs/runtime.png}
\caption{Exact variable-elimination runtime for $P(\text{\var{LUNG\_CANCER}}\mid e)$ versus number of observed variables (mean $\pm$ SD of 30 random queries).}
\label{fig:m3run}
\end{figure}

\section{Revised SOTA Positioning}
\emph{Identify and understand.} The anchor is Neural Graphical Models \cite{ngm}: an MLP trained as multi-task regression under a graph-adherence penalty, with gradient-based inference. Its lung-cancer appendix reports 5-fold CV scores of $0.95\pm0.02$ (LR) and $0.96\pm0.01$ (NGM) for \var{LUNG\_CANCER}, and $0.71\pm0.01$ and $0.79\pm0.02$ for \var{SMOKING}, but does not state the metric.

\emph{Compare and validate.} Our LR score ($0.938\pm0.009$ AUROC) is numerically close to the paper's 0.95, but the metric is unstated and the row count (284 versus 276 de-duplicated, 309 raw) and \var{AGE} handling (binarised versus three bins/continuous) differ; we therefore draw no direct comparison. We did not implement an NGM, so we cannot say whether it would close the gap between our BN and LR; the paper's infant-mortality results favour NGMs over a score-based BN, which is consistent in direction with our finding that a sparse discrete BN trails LR, but this is not evidence about the lung-cancer data.

\emph{Reposition.} M1 promised that a PGM would match standard classifiers while adding interpretability and calibrated uncertainty. The evidence supports a narrower claim:
\begin{itemize}\itemsep0pt
\item \textbf{H1 not supported for discrimination.} The BN trails LR and the SVM by 0.09--0.12 AUROC in all five repeats, and is the worst at finding NO cases at the default threshold. The protected-edge prior does not repair this.
\item \textbf{H2 partly supported.} A stable core of edges exists and agrees with the graphical lasso, but target-adjacent edges are less stable, and the graph grows with $n$ and $\alpha$.
\item \textbf{H3 supported narrowly.} Inference is exact, validated and fast; calibration is best by ECE but not by Brier or log-loss.
\end{itemize}
A PGM is therefore advantageous here for partial-evidence queries, transparent probability computation and inspectable structure, and not for ranking or minority-class detection. M1 items not delivered: hill-climbing/BIC structure learning (replaced by PC), the Naive Bayes baseline, held-out joint log-likelihood (we report target log-loss), the $P(\text{\var{SMOKING}}\mid\text{rest})$ query, a predictive learning curve (sample efficiency is shown for structure only), and the NGM extension. Bootstrap used 100 resamples against 20 planned.

\section{Limitations}
\begin{itemize}\itemsep0pt
\item \textbf{Data:} 276 unique, self-reported, imbalanced rows (38 NO); not a clinical cohort; symptoms may be consequences of disease. Nothing here supports medical use.
\item \textbf{Statistics:} repeated-CV tests are exploratory (shared rows, $n=5$ repeats, Wilcoxon floor 0.0625).
\item \textbf{Baselines untuned:} the MLP's weakness may reflect missing tuning, not model class; thresholds and class weights were not optimised.
\item \textbf{Structure:} no ground truth (no SHD); orientation stability and an alternative CI test for sparse tables were not examined; the protected edges are a prior, not a finding.
\item \textbf{Robustness:} test-time corruption only, single 5-fold split, BN uses imputation rather than marginalisation.
\item \textbf{Bootstrap:} frequencies are read from a figure.
\end{itemize}

\section{Conclusion}
Under repeated, identical, leakage-free evaluation, logistic regression and the SVM are the better predictors of \var{LUNG\_CANCER} in this survey (0.938 and 0.912 AUROC), the PC-learned BN is intermediate (0.822) and the untuned MLP is unreliable (0.566). The BN contributes exact, millisecond-scale partial-evidence inference, a readable graph with a stable core, and good ECE, but not better discrimination, minority-class recall or Brier score. The clinically protected \var{SMOKING} edges are a representational choice with no predictive effect. The learned structure is sample-limited and threshold-dependent, so structural claims should stay tentative. The next experiments are: widen the blanket or add features to the BN, marginalise instead of impute, tune the baselines and decision threshold, test on synthetic graphs with known ground truth (enabling SHD), and train a graph-constrained NGM on the PC graph.

\section*{Reproducibility}
\noindent Python 3.14.2, NumPy 2.4.2, pandas 3.0.3; seed 516; 5 folds $\times$ 5 repeats; 100 bootstrap resamples; PC $\alpha=0.05$, maximum conditioning-set size 3; Dirichlet $\alpha_0=1$. Experiments: \texttt{code\_M3.py} (imports \texttt{Code.py}); outputs in \texttt{results\_m3/}. Figures were drawn from the saved result files; the confidence-interval and paired-test tables were computed from \texttt{repeat\_level\_metrics.csv}.

\msheader{Milestone 4}%
{This report presents the final evaluation of a probabilistic graphical modelling pipeline for the Kaggle Lung Cancer Survey dataset. The study investigates whether a PC-learned Bayesian Network (BN) offers advantages over discriminative classifiers, and identifies the conditions under which its explicit joint distribution and exact inference are useful. After removal of 33 exact duplicate records, the dataset contains 276 observations (238 YES and 38 NO). The principal predictive comparison uses five-fold stratified cross-validation repeated five times, with identical splits for logistic regression, an RBF support-vector machine (SVM), a multilayer perceptron (MLP), and two PC-learned BN variants. Logistic regression obtains the highest AUROC ($0.938\pm0.009$), while the SVM obtains the lowest Brier score ($0.079\pm0.005$) and log-loss ($0.253\pm0.016$). The plain PC-BN obtains AUROC $0.822\pm0.008$ and the lowest expected calibration error (ECE, $0.043\pm0.005$), but has weak minority-class detection at the default threshold (F1 for NO: $0.187\pm0.080$). Four protected \var{SMOKING} edges do not measurably improve predictive performance. Supplementary M4 experiments validate variable elimination against brute-force enumeration for evidence sizes from zero to fifteen, with maximum absolute posterior difference $1.11\times10^{-16}$. On a single stratified 80/20 holdout, exact marginalization of missing variables gives higher AUROC than train-statistic imputation at 10\% and 20\% missingness; these results are exploratory because the test set contains only 56 observations. The evidence does not support the claim that the PGM is the best classifier. Instead, its defensible contribution is an inspectable dependency representation and exact posterior inference under partial evidence.}%
{Probabilistic graphical models, Bayesian networks, PC algorithm, variable elimination, conditional independence, calibration, missing evidence, lung cancer survey.}
\label{ms4}

\section{Introduction}
A discriminative classifier maps a fixed feature vector to a target score or probability. A probabilistic graphical model (PGM), in contrast, represents a joint distribution over variables and can answer conditional probability queries for different evidence subsets. The central question in this project is therefore not whether a PGM wins every predictive metric, but whether its representation and inference capabilities provide useful advantages that conventional classifiers do not directly expose.

The project uses the PC algorithm to estimate a conditional-independence skeleton, orients edges using collider and Meek rules, extends the resulting graph to a DAG for BN parameterization, and estimates conditional probability tables (CPTs) using Dirichlet smoothing. The BN is evaluated alongside logistic regression, an RBF-SVM and an MLP. The M4 analysis consolidates the M1--M3 pipeline and adds explicit inference validation over the full evidence-size range, a held-out joint BN negative log-likelihood (NLL), and a comparison of exact marginalization with imputation under missing evidence.

The study is observational. Neither an undirected dependency nor an arrow in a DAG extension is interpreted as a causal effect.

\section{Problem Formulation}
Let $V=\{X_1,\ldots,X_{15},Y\}$ denote the fifteen survey features and the binary target $Y=\text{\var{LUNG\_CANCER}}$. A BN is a directed acyclic graph $G=(V,E)$ with factorization
\begin{equation}
P(V)=\prod_{i=1}^{|V|}P\big(X_i\mid \mathrm{Pa}(X_i)\big),
\end{equation}
where $\mathrm{Pa}(X_i)$ is the parent set of node $X_i$. For observed evidence $e$, inference computes
\begin{equation}
P(Y=y\mid e)=\frac{P(Y=y,e)}{\sum_{y'}P(Y=y',e)}.
\end{equation}
The project evaluates two distinct capabilities: (1) \emph{prediction}, estimating $P(Y=1\mid X)$ and comparing predictive quality with discriminative baselines; and (2) \emph{partial-evidence inference}, estimating $P(Y\mid e)$ when only a subset of features is observed, marginalizing unobserved variables. These objectives are related but not interchangeable. A model can have lower classification AUROC while still providing a useful general-purpose posterior query interface.

\section{Dataset and Preprocessing}
The source is the Kaggle Lung Cancer Survey dataset \cite{kaggle}. The original file contains 309 rows and 16 columns. Removal of 33 exact duplicate rows leaves 276 unique records: 238 YES and 38 NO. The positive-class proportion is approximately 0.862; consequently, an always-YES classifier has accuracy 0.862. Accuracy alone is therefore not an adequate evaluation measure.

\var{AGE} ranges from 21 to 87. For BN learning, \var{AGE} is discretized into three bins, with cut points fitted using training data in predictive cross-validation. The discriminative pipelines use median imputation and standardization for \var{AGE}; binary features are imputed using the training-fold mode where required. All learned preprocessing transformations are fitted within the training partition. The dataset is self-reported and is not a clinical cohort.

\section{Methods}
\subsection{PC Structure Learning}
The PC algorithm \cite{spirtes,colombo} begins with a complete undirected graph. Candidate edges are tested in increasing conditioning-set size using a conditional likelihood-ratio ($G$) test. An edge is removed when the test does not reject conditional independence at significance level $\alpha=0.05$; the separating set is retained. The maximum conditioning-set size is three. Unshielded colliders are oriented using separating sets, followed by Meek orientation rules. The CPDAG is distinguished from the fully directed DAG extension used for BN inference.

The primary BN uses the plain PC skeleton. An ablation variant restores four edges from \var{SMOKING} to \var{WHEEZING}, \var{YELLOW\_FINGERS}, \var{CHRONIC\_DISEASE} and \var{COUGHING}. These are externally imposed modelling constraints, not edges discovered by PC. Their orientation in the DAG extension is not a causal claim.

\subsection{Parameter Learning and Inference}
CPTs are estimated using Dirichlet smoothing with $\alpha_0=1$. For a node with $r_i$ states, parent configuration $j$, and observed count $N_{ijk}$ for state $k$, the posterior-mean estimate is
\begin{equation}
\hat\theta_{ijk}=\frac{N_{ijk}+\alpha_0}{N_{ij}+r_i\alpha_0}.
\end{equation}
Exact variable elimination (VE) computes posterior queries by multiplying factors, summing out non-query and non-evidence variables, and normalizing the remaining factor. M4 validates VE numerically against brute-force enumeration for evidence sizes from zero through fifteen.

\subsection{Predictive Baselines}
The discriminative baselines are:
\begin{itemize}\itemsep0pt
\item Logistic regression with balanced class weights and maximum 3000 iterations.
\item RBF-SVM with $C=1$, balanced class weights and probability estimates.
\item MLP with hidden layers $(32,16)$, L2 penalty 0.01, early stopping and maximum 1200 iterations.
\end{itemize}
The two BN variants are evaluated on the same folds. No claim is made that the MLP represents the performance of neural networks generally; it is the configured baseline used in this study.

\subsection{Evaluation Protocol and Metrics}
The primary predictive evaluation uses RepeatedStratifiedKFold with five folds and five repeats, random seed 516. Each repeat pools out-of-fold predictions across its five folds, yielding one repeat-level metric observation per model. We report mean, sample standard deviation and 95\% $t$-intervals with four degrees of freedom. Metrics include AUROC, AUPRC, F1 for the NO class at threshold 0.5, Brier score, log-loss, ECE using ten equal-width bins, and accuracy.

\begin{figure}[!htbp]
\centering
\includegraphics[width=\columnwidth]{figs/auroc.png}
\caption{AUROC comparison across the predictive models. Logistic regression ranks highest; the PC-BN is intermediate.}
\label{fig:m4auroc}
\end{figure}

\begin{figure}[!htbp]
\centering
\includegraphics[width=\columnwidth]{figs/f1no.png}
\caption{F1 for the NO class at threshold 0.5. The BN has limited minority-class detection at this operating point.}
\label{fig:m4f1}
\end{figure}

Paired comparisons are computed across the five repeat-level observations. Wilcoxon signed-rank $p$-values are adjusted using Holm's procedure in the recomputed M4 statistical output. These tests remain exploratory: the repeated runs reuse the same 276 records and are not independent external samples \cite{nadeau}. With five pairs, the smallest attainable two-sided Wilcoxon $p$-value is 0.0625.

\section{Predictive Results}
\subsection{Main KPI Comparison}
Table~\ref{tab:m4kpi} reports the principal repeated-CV results. Logistic regression has the highest AUROC, AUPRC and F1 for NO. The SVM has the lowest Brier score and log-loss and the highest accuracy. The plain PC-BN ranks below both LR and SVM on discrimination and probability error, but has the lowest ECE. The protected-edge BN is effectively indistinguishable from the plain BN.

\begin{table}[!htbp]
\centering
\caption{Repeated-CV predictive metrics (mean $\pm$ SD over five repeats of five-fold CV).}
\label{tab:m4kpi}
\tablesize
\setlength{\tabcolsep}{4pt}
\resizebox{\columnwidth}{!}{%
\begin{tabular}{lccccccc}
\toprule
Model & AUROC $\uparrow$ & AUPRC $\uparrow$ & F1(NO) $\uparrow$ & Brier $\downarrow$ & Log-loss $\downarrow$ & ECE $\downarrow$ & Accuracy \\
\midrule
Logistic Regression & 0.938 $\pm$ 0.009 & 0.990 $\pm$ 0.002 & 0.623 $\pm$ 0.027 & 0.091 $\pm$ 0.005 & 0.301 $\pm$ 0.013 & 0.118 $\pm$ 0.003 & 0.864 $\pm$ 0.008 \\
SVM (RBF) & 0.912 $\pm$ 0.020 & 0.984 $\pm$ 0.005 & 0.508 $\pm$ 0.072 & 0.079 $\pm$ 0.005 & 0.253 $\pm$ 0.016 & 0.052 $\pm$ 0.011 & 0.882 $\pm$ 0.011 \\
PC-BN (plain) & 0.822 $\pm$ 0.008 & 0.967 $\pm$ 0.003 & 0.187 $\pm$ 0.080 & 0.099 $\pm$ 0.004 & 0.316 $\pm$ 0.010 & 0.043 $\pm$ 0.005 & 0.854 $\pm$ 0.015 \\
PC-BN (protected) & 0.820 $\pm$ 0.009 & 0.967 $\pm$ 0.003 & 0.170 $\pm$ 0.080 & 0.100 $\pm$ 0.003 & 0.317 $\pm$ 0.010 & 0.040 $\pm$ 0.006 & 0.854 $\pm$ 0.008 \\
MLP & 0.566 $\pm$ 0.127 & 0.883 $\pm$ 0.051 & 0.085 $\pm$ 0.121 & 0.191 $\pm$ 0.021 & 0.572 $\pm$ 0.045 & 0.268 $\pm$ 0.034 & 0.843 $\pm$ 0.015 \\
\bottomrule
\end{tabular}}
\end{table}

\subsection{Probability Quality and Calibration}
The SVM has the best Brier score and log-loss, while the BN has the lowest ECE. These metrics measure different properties: ECE summarizes the average calibration gap across bins, whereas Brier score and log-loss penalize probability errors across individual observations. A low ECE alone does not establish superior probability quality \cite{guo}. Fig.~\ref{fig:m4rel} shows the reliability curves; sparse bins, particularly at the extremes, should be interpreted cautiously because the NO class contains only 38 records.

\begin{figure}[!htbp]
\centering
\includegraphics[width=0.85\columnwidth]{figs/reliab.png}
\caption{Reliability diagram from pooled out-of-fold predictions. The dashed diagonal denotes perfect calibration.}
\label{fig:m4rel}
\end{figure}

\subsection{Error Analysis}
At threshold 0.5, the plain BN recognizes 12.6\% of NO cases, compared with 81.6\% for logistic regression and 44.7\% for SVM. Its sensitivity to YES is 0.971, indicating a strong tendency to predict the majority class. The operating threshold was not optimized. Thus, the F1(NO) result describes the default threshold and should not be interpreted as the best achievable operating point for the BN.

\subsection{Paired Statistical Analysis}
Across the five repeats, LR exceeds the plain BN in AUROC by approximately 0.115 and in F1(NO) by approximately 0.437. SVM exceeds the BN in AUROC by approximately 0.090 and F1(NO) by approximately 0.321. The recomputed M4 paired-test table reports Holm-adjusted Wilcoxon $p$-values; for the LR-versus-BN and SVM-versus-BN comparisons these remain 0.25 for AUROC and 0.25 for F1(NO). The unadjusted Wilcoxon $p$-value is 0.0625, which is the minimum attainable with five paired repeats. Accordingly, the direction and magnitude of differences are more informative than a binary significance claim. The protected-edge ablation shows no measurable predictive improvement.

\section{Representation and Structure-Learning Results}
\subsection{Graph Size and Sensitivity}
For plain PC at $\alpha=0.05$, the learned skeleton contains 18 edges among 16 nodes (density 0.150). The protected variant contains 22 edges. At $\alpha=0.01$, the plain and protected graphs contain 14 and 18 edges; at $\alpha=0.10$, they contain 25 and 26 edges. The target's neighbours remain unchanged between $\alpha=0.01$ and 0.05 and gain \var{ALCOHOL\_CONSUMING} and \var{FATIGUE} at $\alpha=0.10$. This indicates local stability over the lower two thresholds but global sensitivity to the significance level.

\begin{figure}[!htbp]
\centering
\includegraphics[width=\columnwidth]{figs/alpha.png}
\caption{PC skeleton edge count as a function of the significance level, for plain and protected-edge variants.}
\label{fig:m4alpha}
\end{figure}

\subsection{Bootstrap Stability and Sample Size}
The 100-resample bootstrap analysis identifies a stable core of dependencies, while some target-adjacent edges are less stable. The target-neighbour set should therefore not be treated as definitively recovered. The sample-size experiment shows that mean edge count rises from $7.3\pm1.4$ at 55 observations to $16.8\pm1.4$ at 221 observations; the full-data run contains 18 edges. No plateau is visible between 80\% and 100\% of the available data. Subsamples are drawn from the same dataset, so their variability does not represent independent external sampling.

\begin{figure}[!htbp]
\centering
\includegraphics[width=0.85\columnwidth]{figs/bootstrap.png}
\caption{Top bootstrap edge-selection frequencies for the PC skeleton (100 resamples).}
\label{fig:m4boot}
\end{figure}

\begin{figure}[!htbp]
\centering
\includegraphics[width=\columnwidth]{figs/samplesize.png}
\caption{PC skeleton edge count versus training sample size.}
\label{fig:m4size}
\end{figure}

\subsection{Graphical-Lasso Comparison}
The graphical lasso \cite{glasso} returns 82 edges among 16 variables (density 0.683), compared with 18--22 edges for PC. All PC edges are contained in the graphical-lasso graph. The two estimates are not directly interchangeable: graphical lasso estimates a Gaussian conditional-dependence structure, whereas most survey variables are binary. Neither graph is a causal ground truth, and no structural Hamming distance is reported because the real dataset has no known reference graph.

\section{Robustness Under Feature Corruption}
The M3 robustness experiment trains on clean folds and corrupts held-out features using MCAR missingness or feature noise. All models degrade as corruption increases, and the broad ranking LR $\gtrsim$ SVM $>$ BN $>$ MLP is retained. At 20\% MCAR, AUROC is 0.879 for LR, 0.869 for SVM and 0.724 for BN. At 10\% noise, LR and SVM are close (0.884 and 0.882), while BN obtains 0.730. This experiment uses one five-fold split and train-fold imputation; it does not test the BN's exact marginalization capability.

\begin{figure}[!htbp]
\centering
\includegraphics[width=\columnwidth]{figs/robust.png}
\caption{AUROC under held-out feature corruption. These are fold-averaged results from a single five-fold split.}
\label{fig:m4rob}
\end{figure}

\section{M4: Exact Inference Validation}
M4 extends the inference validation across evidence-set sizes from zero to fifteen observed variables. Three random posterior queries are evaluated at each evidence size. VE and brute-force enumeration agree to numerical precision: the maximum absolute posterior difference over the reported cases is $1.11\times10^{-16}$. This supports the numerical correctness of the VE implementation for the tested BN and queries.

The measured mean VE time ranges from approximately 0.16 to 0.56\,ms in this M4 run, while brute-force enumeration takes approximately 2.17 to 3.31\,ms. The timings are implementation- and machine-dependent and should not be generalized as a scalability result. The M4 query count is three per evidence size; it is a separate validation run from the M3 validation, which used 30 queries per size and checked brute force only for evidence sizes zero through four.

\section{M4: Missing-Evidence Marginalization}
A supplementary stratified 80/20 holdout contains 220 training and 56 test observations. The experiment compares exact BN marginalization, in which unobserved features are left unobserved in the posterior query, against train-statistic imputation using identical missingness masks.

At 10\% missingness, marginalization obtains AUROC 0.831, AUPRC 0.964, Brier 0.090 and log-loss 0.304; imputation obtains AUROC 0.725, AUPRC 0.933, Brier 0.102 and log-loss 0.366. At 20\% missingness, marginalization obtains AUROC 0.824, AUPRC 0.965, Brier 0.101 and log-loss 0.326; imputation obtains AUROC 0.574, AUPRC 0.890, Brier 0.127 and log-loss 0.445. In this single split, marginalization is better on all four reported metrics at both missingness levels.

These results are promising but preliminary. The test set contains only 56 observations, the comparison uses one split and one realization of the missingness masks, and the results are not directly comparable to the repeated-CV KPI table. They do not establish statistically significant or generalizable superiority. Repeated holdouts or repeated corruption masks are needed before making a stronger claim.

\begin{table}[!htbp]
\centering
\caption{M4 missing-evidence results on one stratified 80/20 holdout ($n_\mathrm{test}=56$).}
\label{tab:m4missing}
\tablesize
\setlength{\tabcolsep}{4pt}
\begin{tabular}{llcccc}
\toprule
Missing & Method & AUROC & AUPRC & Brier & Log-loss \\
\midrule
10\% & Marginalize & 0.831 & 0.964 & 0.090 & 0.304 \\
10\% & Impute & 0.725 & 0.933 & 0.102 & 0.366 \\
20\% & Marginalize & 0.824 & 0.965 & 0.101 & 0.326 \\
20\% & Impute & 0.574 & 0.890 & 0.127 & 0.445 \\
\bottomrule
\end{tabular}
\end{table}

\section{M4: Joint BN Likelihood}
The supplementary held-out joint likelihood experiment uses a stratified 80/20 split with 220 training observations and 56 test observations. The reported joint negative log-likelihood is 10.3938 nats per test sample (total 582.0523 nats). This is a joint-distribution score over the BN variables and is not the same quantity as target-only classification log-loss in Table~\ref{tab:m4kpi}. It should not be compared numerically with the LR or SVM target log-loss. The result is a single-holdout estimate and is reported as a supplementary characterization of the BN joint model.

\section{Position Relative to Neural Graphical Models}
The selected research anchor is Neural Graphical Models (NGMs) by Shrivastava and Chajewska \cite{ngm}. The cited work combines graph-structured neural parameterizations with graph-adherence regularization and inference procedures. Its lung-cancer appendix reports 5-fold values of $0.95\pm0.02$ for logistic regression and $0.96\pm0.01$ for NGM, but does not state the metric in the relevant table. Its dataset handling also differs from this project. Therefore, these external values are not directly comparable to the present AUROC results.

No NGM was implemented in this project. The present work evaluates a discrete PC-learned BN and conventional discriminative baselines; it does not claim to reproduce or outperform the NGM method. A graph-constrained NGM is a possible future extension, not a completed experiment.

\section{Discussion: Is the PGM Better?}
The results support a task-dependent conclusion:
\begin{itemize}\itemsep0pt
\item \textbf{Predictive discrimination: no.} LR and SVM achieve higher AUROC than the BN. The BN is also weak at recognizing the NO class at the default 0.5 threshold.
\item \textbf{Probability calibration: mixed.} The BN has the lowest ECE, but does not have the best Brier score or log-loss.
\item \textbf{Representation:} the BN provides an explicit dependency graph and a factorized joint distribution. The learned graph has a stable core but remains sensitive to sample size and PC threshold.
\item \textbf{Exact reasoning: yes,} for the tested network and queries. VE matches brute-force enumeration to floating-point precision and supports posterior queries with partial evidence.
\item \textbf{Missing evidence:} preliminary evidence favours exact marginalization over imputation on one holdout. This finding requires repeated validation.
\end{itemize}
Consequently, it is not scientifically justified to state that PGM is universally better than conventional machine learning. The defensible claim is narrower: the PGM sacrifices predictive discrimination in this experiment but provides explicit dependency representation and exact probabilistic reasoning under arbitrary observed evidence subsets.

\section{Limitations}
The main limitations are:
\begin{enumerate}\itemsep0pt
\item The dataset contains only 276 unique self-reported observations, including 38 NO cases, and is not a clinical cohort.
\item Repeated-CV statistical tests reuse the same observations. Their $p$-values are exploratory and do not constitute independent external replication.
\item The M4 missing-evidence and joint-NLL experiments use one 80/20 holdout with 56 test cases.
\item The M4 inference equivalence experiment uses three queries per evidence size; broader randomized validation would provide stronger coverage.
\item PC structure learning depends on the CI test, significance level, conditioning-set cap and finite-sample behavior of sparse contingency tables.
\item The protected \var{SMOKING} edges are imposed modelling constraints, not PC discoveries, and do not improve predictive performance measurably.
\item There is no ground-truth graph for the real survey data; causal effects and structural recovery accuracy cannot be established.
\item The MLP is a compact configured baseline, not an exhaustive neural-model benchmark. No NGM was implemented.
\item No result in this report supports medical diagnosis, screening or treatment decisions.
\end{enumerate}

\section{Conclusion and Future Work}
The final evaluation does not show that the PC-learned BN is a better classifier than logistic regression or RBF-SVM on the Lung Cancer Survey dataset. LR has the strongest discrimination, and SVM has the best Brier score and log-loss. The BN's ECE is lowest, but its minority-class F1 is weak at the default threshold. The protected-edge ablation does not change this outcome.

The PGM's contribution is instead its explicit factorization and ability to answer exact posterior queries when only part of the evidence is observed. M4 verifies VE against brute-force enumeration for evidence sizes zero through fifteen. A single-holdout missing-evidence experiment favours marginalization over imputation, although this finding requires repeated validation.

Future work should repeat missing-evidence experiments over multiple stratified splits and masks, tune decision thresholds and baselines within training data, evaluate alternative CI tests for sparse tables, assess graph stability on larger datasets, and test a graph-constrained NGM. Synthetic data with a known generating graph would enable structural recovery metrics such as structural Hamming distance. External clinical validation would be required before any medical interpretation.

\section*{Reproducibility}
\noindent The experimental configuration uses random seed 516, five-fold stratified cross-validation repeated five times for the main predictive evaluation, PC significance level $\alpha=0.05$, maximum conditioning-set size three, and Dirichlet smoothing $\alpha_0=1$. The M4 supplementary holdout is stratified 80/20 (220 train, 56 test). The final code and result-generation commands should be archived with the submission. Reported M4 holdout results must be regenerated if the code, environment or randomization protocol changes.

\label{refs}
\begin{thebibliography}{9}\small
\bibitem{ngm} H. Shrivastava and U. Chajewska, ``Neural Graphical Models,'' arXiv:2210.00453, 2022. [Online]. Available: \url{https://arxiv.org/abs/2210.00453v4}
\bibitem{kaggle} N. Alaswad, ``Lung Cancer Dataset,'' Kaggle. [Online]. Available: \url{https://www.kaggle.com/datasets/nancyalaswad90/lung-cancer}
\bibitem{spirtes} P. Spirtes, C. Glymour, and R. Scheines, \emph{Causation, Prediction, and Search}, 2nd ed. Cambridge, MA, USA: MIT Press, 2000.
\bibitem{colombo} D. Colombo and M. H. Maathuis, ``Order-independent constraint-based causal structure learning,'' \emph{J. Mach. Learn. Res.}, vol. 15, pp. 3741--3782, 2014.
\bibitem{glasso} J. Friedman, T. Hastie, and R. Tibshirani, ``Sparse inverse covariance estimation with the graphical lasso,'' \emph{Biostatistics}, vol. 9, no. 3, pp. 432--441, 2008.
\bibitem{guo} C. Guo, G. Pleiss, Y. Sun, and K. Q. Weinberger, ``On calibration of modern neural networks,'' in \emph{Proc. ICML}, 2017, pp. 1321--1330.
\bibitem{nadeau} C. Nadeau and Y. Bengio, ``Inference for the generalization error,'' \emph{Mach. Learn.}, vol. 52, pp. 239--281, 2003.
\bibitem{cdc} ``Data Access -- Vital Statistics Online,'' 2019. [Online]. Available: \url{https://www.cdc.gov/nchs/data_access/vitalstatsonline.htm}
\bibitem{ngmcode} H. Shrivastava, ``neural-graphical-models,'' GitHub. [Online]. Available: \url{https://github.com/harshs27/neural-graphical-models}
\end{thebibliography}
\end{document}
